% ============================================================================
%  CHEM 5200: Computational Chemistry
%  Companion notes for Computational Set 2 -- Two-Level Electronic Systems
%
%  Compile with pdfLaTeX (twice, for cross-references).
%  On Overleaf: Menu -> Compiler -> pdfLaTeX.
%
%  To produce the STUDENT version (problems without solutions), change
%  \solutionstrue to \solutionsfalse in the line below.
% ============================================================================
\documentclass[11pt]{article}
\usepackage{preamble}

\solutionstrue      % <-- \solutionsfalse for the student handout

\title{\vspace{-1.2cm}\Large\bfseries Two-Level Electronic Systems\\[2pt]
       \large\normalfont Companion notes to Computational Set 2}
\author{\normalsize Jay Foley\\[-2pt]
        \small Department of Chemistry, University of North Carolina Charlotte}
\date{\small CHEM 5200 \textperiodcentered\ Computational Chemistry}

\begin{document}
\maketitle
\thispagestyle{fancy}

\begin{abstract}
\noindent
\small
In Computational Set 1 we developed the formalism of spin~$1/2$: the Pauli matrices,
their eigenvectors, and their commutation relations. Nothing in that development
required the system to be a spin. These notes make the general statement precise ---
\emph{every} quantum system with two relevant states is governed by the same algebra ---
and then apply it to four concrete chemical problems: the two-orbital LCAO dimer, a
spin or chromophore driven by an oscillating field, donor--acceptor electron transfer
along a nuclear coordinate, and the ammonia inversion doublet in a static electric
field. In each case a $\szo$ term measures the energy asymmetry between the two basis
states and a $\sxo$ (or $\syo$) term measures the coupling between them. We close by
asking where a genuine $\syo$ contribution can come from, and by identifying the two
models --- Jaynes--Cummings and Holstein --- that these examples become once a boson
mode is added.
\end{abstract}

\vspace{-0.3cm}
\begin{nbbox}[How to use these notes]
These notes are the analytical companion to the Jupyter notebook
\texttt{notebooks/two\_level\_systems.ipynb}. The notebook builds the same results
numerically and asks you to write the code; the notes carry the derivations that
are only sketched there. Read \cref{sec:formalism,sec:geometry} before starting the
notebook, and use the appendices whenever a step in the main text is quoted
rather than shown. The problems in \cref{sec:problems} are the same ones that close
the notebook.
\end{nbbox}

\vspace{0.2cm}
{\small\tableofcontents}
\clearpage

% ============================================================================
\section{Every two-level Hamiltonian is a spin}
\label{sec:formalism}
% ============================================================================

Chemistry supplies two-level systems in abundance. A bond in a minimal basis has two
atomic orbitals; a proton in an NMR magnet has two spin projections; a donor--acceptor
complex has the electron on one site or the other; ammonia has its umbrella up or down;
a chromophore in a weak laser field has a ground and a first excited state. In each case
the relevant Hilbert space is two dimensional, and in each case --- as we now show ---
the Hamiltonian is a linear combination of the identity and the three Pauli matrices.

\subsection{The Pauli matrices are a complete basis}

Let $\ket{\alpha}$ and $\ket{\beta}$ be an orthonormal pair spanning the space. Which
two physical states these are will change from section to section; the algebra will not.
Any Hamiltonian is a $2\times2$ Hermitian matrix,
%
\begin{equation}
\Hop =
\begin{bmatrix}
H_{\alpha\alpha} & H_{\alpha\beta}\\
H_{\alpha\beta}^{*} & H_{\beta\beta}
\end{bmatrix},
\qquad H_{\alpha\alpha}, H_{\beta\beta}\in\mathbb{R},
\label{eq:generalH}
\end{equation}
%
which carries four real degrees of freedom: the two diagonal entries and the real and
imaginary parts of $H_{\alpha\beta}$. The set $\{\Iop,\sxo,\syo,\szo\}$ carries four
real degrees of freedom as well, and it is straightforward to check that these four
matrices are linearly independent. They therefore span the space of Hermitian
$2\times2$ matrices, and we may always write

\begin{keyresult}[The Pauli decomposition]
\vspace{-0.35cm}
\begin{equation}
\Hop \;=\; \varepsilon_0\,\Iop \;+\; \half\!\left(d_x\sxo + d_y\syo + d_z\szo\right)
    \;\equiv\; \varepsilon_0\,\Iop + \half\,\dvec\cdot\pauliv ,
\label{eq:pauli-decomp}
\end{equation}
with $\varepsilon_0$ and the three components of $\dvec$ all real.
\end{keyresult}

\noindent
Writing \cref{eq:pauli-decomp} out in components,
%
\begin{equation}
\Hop =
\begin{bmatrix}
\varepsilon_0 + \tfrac{d_z}{2} & \tfrac{d_x - \ii d_y}{2}\\[4pt]
\tfrac{d_x + \ii d_y}{2} & \varepsilon_0 - \tfrac{d_z}{2}
\end{bmatrix},
\label{eq:pauli-components}
\end{equation}
%
so that each piece can be read off and given a name.

\begin{center}
\small
\begin{tabular}{@{}lll@{}}
\toprule
Component & Definition & Physical meaning \\
\midrule
$\varepsilon_0$ & $\half(H_{\alpha\alpha}+H_{\beta\beta})$
  & mean energy; shifts everything, changes nothing \\
$d_z$ & $H_{\alpha\alpha}-H_{\beta\beta}$
  & \textbf{asymmetry}: the energy gap between the basis states \\
$d_x$ & $2\,\mathrm{Re}\,H_{\alpha\beta}$
  & \textbf{coupling}: the real part of the mixing element \\
$d_y$ & $-2\,\mathrm{Im}\,H_{\alpha\beta}$
  & \textbf{coupling}: the imaginary part \\
\bottomrule
\end{tabular}
\end{center}

That table is the thesis of these notes. Everything that follows is an exercise in
identifying, for some piece of chemistry, what plays the role of the asymmetry and what
plays the role of the coupling.

\subsection{Extracting the components}

The Pauli matrices satisfy the two trace identities
%
\begin{equation}
\Tr\!\left(\sxo\right)=\Tr\!\left(\syo\right)=\Tr\!\left(\szo\right)=0,
\qquad
\Tr\!\left(\hat\sigma_i\hat\sigma_j\right) = 2\,\delta_{ij},
\label{eq:traces}
\end{equation}
%
which follow immediately from the product rule
$\hat\sigma_i\hat\sigma_j = \delta_{ij}\Iop + \ii\,\epsilon_{ijk}\hat\sigma_k$
(derived in \cref{app:pauli}). Multiplying \cref{eq:pauli-decomp} by $\hat\sigma_i$ and
tracing therefore isolates a single component:

\begin{equation}
\boxed{\;\varepsilon_0 = \half\Tr\!\left(\Hop\right),
\qquad
d_i = \Tr\!\left(\hat\sigma_i\,\Hop\right).\;}
\label{eq:extract}
\end{equation}

\begin{nbbox}[In the notebook]
\Cref{eq:extract} is the function \texttt{pauli\_decomposition(H)}, and
\cref{eq:pauli-decomp} is \texttt{build\_two\_level(e0, dz, dx, dy)}. The notebook
verifies that the two are exact inverses of one another on 500 randomly generated
Hermitian matrices --- a much stronger check than any single worked example, because
it tests completeness of the basis rather than one lucky case.
\end{nbbox}

\subsection{A first example: the spin operators themselves}

As a sanity check, take the operators from Computational Set 1. Since
$\Szo = \tfrac{\hbar}{2}\szo$, applying \cref{eq:extract} gives
$\varepsilon_0=0$ and $\dvec = (0,0,\hbar)$. Similarly $\Sxo$ gives
$\dvec=(\hbar,0,0)$ and $\Syo$ gives $\dvec=(0,\hbar,0)$. The three spin operators are
simply the three coordinate axes of the $\dvec$ space, which is why that space is worth
taking seriously as a geometric object --- the subject of the next section.

% ============================================================================
\section{Solving the general two-level problem}
\label{sec:geometry}
% ============================================================================

\subsection{Eigenvalues without diagonalizing}

The decomposition of \cref{sec:formalism} pays for itself immediately. Using the
product rule for Pauli matrices one finds (\cref{app:pauli})
%
\begin{equation}
\left(\dvec\cdot\pauliv\right)^2 = |\dvec|^2\,\Iop .
\label{eq:dsquared}
\end{equation}
%
A matrix whose square is a multiple of the identity can only have eigenvalues
$\pm|\dvec|$, so the spectrum of $\Hop$ follows with no linear algebra at all:

\begin{keyresult}[Splitting]
\vspace{-0.35cm}
\begin{equation}
E_{\pm} = \varepsilon_0 \pm \frac{|\dvec|}{2}
        = \varepsilon_0 \pm \half\sqrt{d_x^2+d_y^2+d_z^2},
\qquad
E_+ - E_- = |\dvec| .
\label{eq:eigenvalues}
\end{equation}
\end{keyresult}

\noindent
This one formula is the bonding--antibonding splitting, the generalized Rabi frequency,
the adiabatic gap at an avoided crossing, the Davydov splitting of an excitonic dimer,
and the inversion doublet of ammonia. It is worth pausing on what it says. Writing
$\dperp \equiv \sqrt{d_x^2+d_y^2}$ for the coupling magnitude,
%
\begin{equation}
E_+ - E_- = \sqrt{d_z^2 + \dperp^2} \;\geq\; \dperp .
\label{eq:noncrossing}
\end{equation}
%
The asymmetry and the coupling add \emph{in quadrature}, so the splitting can never fall
below the coupling. Even when the two basis states are exactly degenerate, they are
split by $\dperp$. This is the \textbf{non-crossing rule}: two states that are coupled
cannot become degenerate by tuning one parameter.

\subsection{The mixing angle}

Write $\dvec$ in spherical polar coordinates,
%
\begin{equation}
d_z = |\dvec|\cos\theta,\qquad
d_x = |\dvec|\sin\theta\cos\varphi,\qquad
d_y = |\dvec|\sin\theta\sin\varphi,
\end{equation}
%
so that
%
\begin{equation}
\tan\theta = \frac{\dperp}{d_z},
\qquad
\varphi = \atantwo(d_y,d_x).
\label{eq:mixing-angle}
\end{equation}
%
Substituting into \cref{eq:pauli-components} and solving (\cref{app:eigvec}) gives the
eigenvectors:

\begin{keyresult}[Eigenvectors]
\vspace{-0.35cm}
\begin{align}
\ket{+} &= \cos\tfrac{\theta}{2}\,\ket{\alpha} + e^{\ii\varphi}\sin\tfrac{\theta}{2}\,\ket{\beta},\\
\ket{-} &= -e^{-\ii\varphi}\sin\tfrac{\theta}{2}\,\ket{\alpha} + \cos\tfrac{\theta}{2}\,\ket{\beta}.
\label{eq:eigenvectors}
\end{align}
\end{keyresult}

\noindent
The angle $\theta$ is the single most useful number in two-level physics: it measures how
thoroughly the coupling has scrambled the basis states. The population of $\ket{\alpha}$
in the upper state is $\cos^2(\theta/2)$ and in the lower state $\sin^2(\theta/2)$.

\begin{center}
\small
\begin{tabular}{@{}lll@{}}
\toprule
Regime & $\theta$ & Eigenstates are\ldots \\
\midrule
$d_z \gg \dperp > 0$   & $\to 0$     & essentially $\ket{\alpha},\ket{\beta}$; coupling is a perturbation \\
$d_z = 0$              & $=\pi/2$    & exactly $50/50$; maximally delocalized \\
$-d_z \gg \dperp > 0$  & $\to \pi$   & localized again, with the labels exchanged \\
\bottomrule
\end{tabular}
\end{center}

\noindent
Note the half-angles in \cref{eq:eigenvectors}. As $d_z$ sweeps from $+\infty$ to
$-\infty$ at fixed $\dperp$, $\theta$ runs from $0$ to $\pi$ and $\ket{+}$ evolves
continuously from $\ket{\alpha}$ to $\ket{\beta}$. That continuous exchange of character,
without the levels ever touching, \emph{is} an avoided crossing. \Cref{fig:avoided}
shows both halves of the story.

% ---------------------------------------------------------------------------
\begin{figure}[t]
\centering
\begin{tikzpicture}
\begin{groupplot}[
  group style={group size=2 by 1, horizontal sep=1.6cm},
  notestyle, height=6.2cm,
  legend style={at={(0.5,-0.30)}, anchor=north, legend columns=2,
                draw=none, fill=none, font=\footnotesize},
  legend cell align=left,
]
% ---- left: the avoided crossing
\nextgroupplot[
  xlabel={$d_z$}, ylabel={energy},
  title={Levels repel},
  xmin=-6, xmax=6, ymin=-3.4, ymax=3.4,
]
\addplot[cMuted, dashed, thin, domain=-6:6, samples=2]{0.5*x};
\addlegendentry{uncoupled $H_{\alpha\alpha}$}
\addplot[cMuted, dotted, thin, domain=-6:6, samples=2]{-0.5*x};
\addlegendentry{uncoupled $H_{\beta\beta}$}
\addplot[cUpper, very thick, domain=-6:6, samples=200]{0.5*sqrt(x^2+1)};
\addlegendentry{$E_+$}
\addplot[cLower, very thick, domain=-6:6, samples=200]{-0.5*sqrt(x^2+1)};
\addlegendentry{$E_-$}
\draw[<->, cAccent, thick] (axis cs:0,-0.5) -- (axis cs:0,0.5);
\node[cAccent, anchor=west, font=\footnotesize] at (axis cs:0.35,0){$\dperp$};

% ---- right: the composition
\nextgroupplot[
  xlabel={$d_z$}, ylabel={population in $\ket{-}$},
  title={Character is exchanged},
  xmin=-6, xmax=6, ymin=-0.04, ymax=1.04,
]
\addplot[cLower, very thick, domain=-6:6, samples=300]
  {sin(0.5*atan2(1,x))^2};
\addlegendentry{$|\!\braket{\alpha|-}\!|^2$}
\addplot[cUpper, very thick, domain=-6:6, samples=300]
  {cos(0.5*atan2(1,x))^2};
\addlegendentry{$|\!\braket{\beta|-}\!|^2$}
\addplot[cMuted, dotted, thin, domain=-6:6, samples=2]{0.5};
\end{groupplot}
\end{tikzpicture}
\caption{The two faces of an avoided crossing, for fixed coupling $\dperp = 1$.
\textbf{Left:} the uncoupled diagonal energies (grey) cross at $d_z=0$; the
eigenvalues of \cref{eq:eigenvalues} do not, and are separated there by exactly
$\dperp$. \textbf{Right:} the composition of the lower eigenstate, from
\cref{eq:eigenvectors}. The exchange of character happens over a window of width
$\sim\dperp$ in $d_z$ --- a stronger coupling makes the switch \emph{more} gradual,
not less.}
\label{fig:avoided}
\end{figure}
% ---------------------------------------------------------------------------

\subsection{Dynamics: the Bloch equation}
\label{sec:bloch}

For a normalized state $\ket{\psi}$ define the \textbf{Bloch vector}
%
\begin{equation}
\svec = \Braket{\pauliv}
      = \Big(\Braket{\psi|\sxo|\psi},\;\Braket{\psi|\syo|\psi},\;\Braket{\psi|\szo|\psi}\Big).
\label{eq:bloch-vector}
\end{equation}
%
Three real numbers, and for any pure state $|\svec| = 1$ exactly, so the state lives on
the surface of a unit sphere. The north pole is $\ket{\alpha}$, the south pole is
$\ket{\beta}$, the equator holds the equally weighted superpositions, and longitude on
the equator is the relative phase. Comparing \cref{eq:bloch-vector} with
\cref{eq:eigenvectors}, the eigenstate $\ket{+}$ sits at polar angles $(\theta,\varphi)$
--- that is, $\svec_+ = \dhat$. \textbf{The eigenstates are the two points where the
state vector is parallel or antiparallel to $\dvec$.}

Applying the Heisenberg equation to each Pauli matrix and using
$[\hat\sigma_a,\hat\sigma_b]=2\ii\,\epsilon_{abc}\hat\sigma_c$ gives a strikingly simple
result, worked through in full in \cref{app:blocheq}:

\begin{keyresult}[Bloch equation]
\vspace{-0.35cm}
\begin{equation}
\frac{\dd\svec}{\dd t} = \frac{1}{\hbar}\,\dvec\times\svec .
\label{eq:bloch-eq}
\end{equation}
\end{keyresult}

\noindent
The Bloch vector \emph{precesses about $\dvec$} at angular frequency $|\dvec|/\hbar$.
Three consequences deserve to be stated explicitly.

\begin{enumerate}[leftmargin=*,itemsep=2pt]
\item \textbf{Energy is conserved.} A cross product is perpendicular to both arguments,
so $\dvec\cdot\svec$ is constant --- which is exactly the statement that
$\Braket{\Hop} = \varepsilon_0 + \half\dvec\cdot\svec$ does not change.

\item \textbf{Eigenstates do not move.} If $\svec \parallel \dvec$ the cross product
vanishes. Stationary states are stationary.

\item \textbf{The precession frequency is the Bohr frequency.} Since
$|\dvec| = E_+-E_-$, the rate is $(E_+-E_-)/\hbar$.
\end{enumerate}

\noindent
Notice that $\varepsilon_0$ is absent from \cref{eq:bloch-eq}. An overall energy shift
produces an overall phase, which no measurement on this system can detect.

% ---------------------------------------------------------------------------
\begin{figure}[t]
\centering
\tdplotsetmaincoords{72}{118}
\begin{subfigure}[b]{0.47\textwidth}
\centering
\begin{tikzpicture}[scale=2.05]
  % --- screen-space silhouette (NOT in 3d coords, so it stays a circle) ---
  \shade[ball color=gray!22, opacity=0.16] (0,0) circle (1);
  \draw[gray!50, thin] (0,0) circle (1);
  \begin{scope}[tdplot_main_coords]
    % equator and the xz meridian
    \tdplotdrawarc[gray!45,thin,densely dashed]{(0,0,0)}{1}{0}{360}{}{}
    \tdplotsetrotatedcoords{0}{90}{0}
    \tdplotdrawarc[tdplot_rotated_coords,gray!35,thin,densely dotted]{(0,0,0)}{1}{0}{360}{}{}
    \tdplotsetrotatedcoords{0}{0}{0}
    % axes
    \draw[-{Stealth[length=4pt]}, gray!65] (0,0,0) -- (1.38,0,0)
      node[anchor=north,black!65,font=\scriptsize]{$x$};
    \draw[-{Stealth[length=4pt]}, gray!65] (0,0,0) -- (0,1.38,0)
      node[anchor=west,black!65,font=\scriptsize]{$y$};
    \draw[-{Stealth[length=4pt]}, gray!65] (0,0,0) -- (0,0,1.30)
      node[anchor=south,black!65,font=\scriptsize]{$z$};
    % the d vector at theta = 54, phi = 40
    \tdplotsetcoord{P}{1}{54}{40}
    \tdplotsetcoord{Pxy}{0.809}{90}{40}
    \draw[gray!65, dashed, thin] (P) -- (Pxy);
    \draw[gray!65, dashed, thin] (0,0,0) -- (Pxy);
    \draw[-{Stealth[length=6pt]}, cAccent, very thick] (0,0,0) -- (P);
    \node[cAccent, anchor=south west, font=\small, inner sep=1pt] at (P) {$\dhat$};
    % theta arc, in the plane containing z and d
    \tdplotsetthetaplanecoords{40}
    \tdplotdrawarc[tdplot_rotated_coords,cAccent,thin,-{Stealth[length=3.5pt]}]%
      {(0,0,0)}{0.30}{0}{54}{anchor=east,font=\scriptsize,cAccent,inner sep=1.5pt}{$\theta$}
    % phi arc, in the xy plane
    \tdplotdrawarc[cUpper,thin,-{Stealth[length=3.5pt]}]{(0,0,0)}{0.52}{0}{40}%
      {anchor=north west,font=\scriptsize,cUpper,inner sep=1pt}{$\varphi$}
    % poles
    \fill[black] (0,0,1) circle (0.9pt);
    \node[anchor=east, font=\scriptsize, inner sep=2pt] at (0,0,1) {$\ket{\alpha}$};
    \fill[black] (0,0,-1) circle (0.9pt);
    \node[anchor=east, font=\scriptsize, inner sep=2pt] at (0,0,-1) {$\ket{\beta}$};
  \end{scope}
\end{tikzpicture}
\caption{Geometry. $\dvec$ has polar angles $(\theta,\varphi)$; the eigenstate
$\ket{+}$ points along $\dhat$ and $\ket{-}$ along $-\dhat$.}
\label{fig:bloch-geom}
\end{subfigure}
\hfill
\begin{subfigure}[b]{0.47\textwidth}
\centering
\begin{tikzpicture}[scale=2.05]
  \shade[ball color=gray!22, opacity=0.16] (0,0) circle (1);
  \draw[gray!50, thin] (0,0) circle (1);
  \begin{scope}[tdplot_main_coords]
    \tdplotdrawarc[gray!45,thin,densely dashed]{(0,0,0)}{1}{0}{360}{}{}
    \draw[-{Stealth[length=4pt]}, gray!65] (0,0,0) -- (1.38,0,0)
      node[anchor=north,black!65,font=\scriptsize]{$x$};
    \draw[-{Stealth[length=4pt]}, gray!65] (0,0,0) -- (0,1.38,0)
      node[anchor=west,black!65,font=\scriptsize]{$y$};
    \draw[-{Stealth[length=4pt]}, gray!65] (0,0,0) -- (0,0,1.30)
      node[anchor=south,black!65,font=\scriptsize]{$z$};
    % precession cone: half-angle 54.5 deg about d-hat, which lies in the xz plane
    \tdplotsetrotatedcoords{0}{54.5}{0}
    \tdplotdrawarc[tdplot_rotated_coords, cUpper, very thick]%
       {(0,0,0.5807)}{0.8141}{0}{360}{}{}
    \tdplotsetrotatedcoords{0}{0}{0}
    % the precession axis
    \tdplotsetcoord{D}{1.16}{54.5}{0}
    \draw[-{Stealth[length=6pt]}, cAccent, very thick] (0,0,0) -- (D);
    \node[cAccent, anchor=south east, font=\small, inner sep=3pt] at (D) {$\dhat$};
    % start of the trajectory, and a quarter period later
    % (s(pi/2) = s_par + d-hat x s_perp, worked out for a 54.5 deg cone)
    \fill[black] (0,0,1) circle (1.3pt);
    \fill[cUpper] (0.4728,-0.8141,0.3372) circle (1.3pt);
    \node[anchor=south, font=\scriptsize, inner sep=2pt] at (0,0,1.02) {$\svec(0)=\ket{\alpha}$};
    \fill[black] (0,0,-1) circle (0.9pt);
    \node[anchor=east, font=\scriptsize, inner sep=2pt] at (0,0,-1) {$\ket{\beta}$};
  \end{scope}
\end{tikzpicture}
\caption{Dynamics. Starting at $\ket{\alpha}$, the state sweeps a cone about $\dhat$
at rate $|\dvec|/\hbar$ (red dot: a quarter period later).}
\label{fig:bloch-prec}
\end{subfigure}
\caption{The Bloch sphere. Because the trajectory is a cone about $\dhat$, a state
prepared at the north pole reaches the south pole only if $\dhat$ lies in the
equatorial plane --- which is precisely the resonance condition of \cref{sec:driven}.}
\label{fig:bloch}
\end{figure}
% ---------------------------------------------------------------------------

\begin{nbbox}[In the notebook]
\Cref{eq:bloch-vector} is \texttt{bloch\_vector(psi)}; time evolution is
\texttt{propagate(H, psi0, times)}, which phases the eigenbasis coefficients rather
than integrating an ODE. The notebook checks \cref{eq:bloch-eq} indirectly, by
confirming that the projection $\svec\cdot\dhat$ is constant to machine precision
along a trajectory.
\end{nbbox}

% ============================================================================
\section{Example I: the two-orbital LCAO dimer}
\label{sec:dimer}
% ============================================================================

The simplest molecular orbital problem in chemistry is two atoms carrying one atomic
orbital each. Identify $\ket{A}\equiv\ket{\alpha}$ and $\ket{B}\equiv\ket{\beta}$.

\begin{nbbox}[A word on orthogonality]
Atomic orbitals on neighbouring atoms overlap, $\braket{A|B}=S\neq 0$, which strictly
gives a generalized eigenvalue problem $\mathbf{Hc} = E\,\mathbf{Sc}$. Throughout this
section we assume the basis has already been symmetrically (L\"owdin) orthogonalized,
$\mathbf{S}^{-1/2}\mathbf{H}\mathbf{S}^{-1/2}$, so that $\mathbf{S}=\mathbf{1}$. This is
exactly the assumption H\"uckel theory makes, and it is what lets us use the machinery
of \cref{sec:geometry} unmodified. The general overlap case is taken up in the Basis
Sets set.
\end{nbbox}

\subsection{Identifying $\dvec$}

In the orthogonalized basis the matrix elements carry their traditional H\"uckel names,
%
\begin{equation}
\Hop = \begin{bmatrix} \alpha_A & \beta \\ \beta & \alpha_B \end{bmatrix},
\end{equation}
%
where $\alpha_X = \braket{X|\Hop|X}$ is the \emph{Coulomb integral} --- the energy of an
electron localized in the orbital on atom $X$, more negative for a more electronegative
atom --- and $\beta = \braket{A|\Hop|B}$ is the \emph{resonance} or \emph{hopping
integral}, negative by convention and falling off roughly exponentially with
internuclear separation. Applying \cref{eq:extract},

\begin{keyresult}[LCAO dimer]
\vspace{-0.35cm}
\begin{equation}
\varepsilon_0 = \frac{\alpha_A+\alpha_B}{2},
\qquad
d_z = \underbrace{\alpha_A-\alpha_B}_{\text{electronegativity difference}},
\qquad
d_x = \underbrace{2\beta}_{\text{twice the hopping}},
\qquad
d_y = 0 .
\label{eq:dimer-d}
\end{equation}
\end{keyresult}

\noindent
The absence of $d_y$ is not an accident of this example; it is a consequence of
time-reversal symmetry, taken up in \cref{sec:sigmay}.

From \cref{eq:eigenvalues} the molecular orbital energies are
%
\begin{equation}
E_{\pm} = \frac{\alpha_A+\alpha_B}{2} \pm \sqrt{\left(\frac{\alpha_A-\alpha_B}{2}\right)^{\!2} + \beta^2},
\label{eq:dimer-E}
\end{equation}
%
and the mixing angle is $\tan\theta = 2\beta/(\alpha_A-\alpha_B)$.

\subsection{The two limits}

\paragraph{Homonuclear: $d_z=0$.} For H$_2^+$, or for the $\pi$ system of ethylene,
$\alpha_A=\alpha_B\equiv\alpha$ and the Hamiltonian is \emph{pure} $\sxo$. Then
$\theta=\pi/2$, the eigenvectors are $\tfrac{1}{\sqrt2}(\ket{A}\pm\ket{B})$, and
$E_\pm = \alpha \pm |\beta|$: a bonding orbital at $\alpha+\beta$ (recall $\beta<0$) and
an antibonding orbital at $\alpha-\beta$, split by $2|\beta|$.

It is worth noticing that this Hamiltonian is $\Sxo$ up to a scale and a shift. The
bonding and antibonding MOs of a homonuclear diatomic are, as vectors, precisely the
eigenvectors of $\Sxo$ you computed in Computational Set 1.

\paragraph{Strongly heteronuclear: $|d_z| \gg |d_x|$.} Now $\theta\to 0$, the orbitals
barely mix, and the doubly occupied lower MO sits almost entirely on the more
electronegative atom. This is the ionic limit. Everything in between is a polar
covalent bond, and the single dimensionless parameter controlling where you sit is
%
\begin{equation}
\frac{|d_x|}{|d_z|} = \frac{2|\beta|}{|\alpha_A-\alpha_B|}.
\label{eq:covalency}
\end{equation}

\begin{nbbox}[Mixed-valence chemistry]
Ratio \eqref{eq:covalency} is the physical content of the \textbf{Robin--Day}
classification of mixed-valence complexes \cite{RobinDay}. Class~I ($|\beta|\ll|d_z|/2$)
means fully trapped valences and no intervalence band; Class~III
($|\beta|\gg|d_z|/2$) means a completely delocalized, symmetric ion --- the
Creutz--Taube complex
$[(\mathrm{NH_3})_5\mathrm{Ru}\text{-pyrazine-}\mathrm{Ru(NH_3)_5}]^{5+}$
being the canonical case \cite{CreutzTaube}. Class~II is the interesting middle, where
the intervalence charge-transfer band appears and Marcus theory (\cref{sec:et})
governs the thermal transfer rate.
\end{nbbox}

\subsection{A worked case}

Take $\alpha_A = \SI{-11.4}{\electronvolt}$, $\alpha_B = \SI{-13.6}{\electronvolt}$
(atom B the more electronegative), and $\beta = \SI{-2.7}{\electronvolt}$. Then
$\varepsilon_0 = \SI{-12.5}{\electronvolt}$, $d_z = \SI{2.2}{\electronvolt}$, and
$d_x = \SI{-5.4}{\electronvolt}$, so
%
\begin{equation}
\dperp = |2\beta| = \SI{5.4}{\electronvolt},
\qquad
|\dvec| = \sqrt{2.2^2+5.4^2}\;\si{\electronvolt} = \SI{5.831}{\electronvolt},
\end{equation}
%
and, since $d_x<0$ puts the transverse component along $-\hat{x}$,
%
\begin{equation}
\theta = \atantwo(\dperp,\,d_z) = \atantwo(5.4,\,2.2) = \ang{67.83},
\qquad
\varphi = \pi .
\end{equation}
%
The MOs come out at $\SI{-15.415}{\electronvolt}$ (bonding) and
$\SI{-9.585}{\electronvolt}$ (antibonding). The bonding MO carries
$\sin^2(\theta/2) = 0.311$ of its density on A and $0.689$ on B --- polarized toward the
electronegative atom, as chemical intuition demands, but nowhere near the ionic limit
because $2|\beta|$ is more than twice $|d_z|$.

\Cref{fig:avoided} is exactly the picture of this system as $d_z$ is scanned at fixed
$\beta$: a bond that is ionic at one extreme, purely covalent at $d_z=0$, and ionic with
the opposite polarity at the other extreme.

% ---------------------------------------------------------------------------
\begin{figure}[t]
\centering
\begin{tikzpicture}[x=1cm,y=1cm,
    lvl/.style={line width=1.2pt},
    corr/.style={gray!55,line width=0.4pt},
    lab/.style={font=\scriptsize},
    elec/.style={-{Stealth[length=3.2pt]},black,line width=0.5pt}]

  % ================= homonuclear =================
  \begin{scope}[xshift=0cm]
    \node[font=\small\bfseries] at (1.9,2.55) {Homonuclear: $d_z=0$};
    % AO levels
    \draw[cMuted,lvl] (0.0,0) -- (0.8,0);
    \draw[cMuted,lvl] (3.0,0) -- (3.8,0);
    \node[lab,anchor=south] at (0.4,0.10) {$\alpha$};
    \node[lab,anchor=south] at (3.4,0.10) {$\alpha$};
    % MO levels
    \draw[cUpper,lvl] (1.5,1.15) -- (2.3,1.15);
    \draw[cLower,lvl] (1.5,-1.15) -- (2.3,-1.15);
    \node[lab,anchor=south] at (1.9,1.24) {$\alpha-\beta$ \; (antibonding)};
    \node[lab,anchor=north] at (1.9,-1.55) {$\alpha+\beta$ \; (bonding)};
    % correlation lines
    \foreach \sx in {0.8,3.0} {
      \draw[corr] (\sx,0) -- (1.9,1.15);
      \draw[corr] (\sx,0) -- (1.9,-1.15);
    }
    % electron pair in the bonding MO
    \draw[elec] (1.82,-1.35) -- (1.82,-0.95);
    \draw[elec] (1.98,-0.95) -- (1.98,-1.35);
    % splitting arrow, clear to the right of everything
    \draw[<->,cAccent,line width=0.7pt] (4.55,-1.15) -- (4.55,1.15);
    \node[cAccent,lab,anchor=west] at (4.65,0.0) {$2|\beta|$};
    \node[lab,anchor=north,align=center] at (1.9,-2.10)
      {$50/50$ on both atoms\\[-2pt]($\theta=\pi/2$)};
  \end{scope}

  % ================= heteronuclear =================
  \begin{scope}[xshift=7.7cm]
    \node[font=\small\bfseries] at (1.9,2.55) {Heteronuclear: $d_z\neq0$};
    \draw[cMuted,lvl] (0.0,0.75) -- (0.8,0.75);
    \draw[cMuted,lvl] (3.0,-0.75) -- (3.8,-0.75);
    \node[lab,anchor=south] at (0.4,0.85) {$\alpha_A$};
    \node[lab,anchor=south] at (3.4,-0.65) {$\alpha_B$};
    \draw[cUpper,lvl] (1.5,1.55) -- (2.3,1.55);
    \draw[cLower,lvl] (1.5,-1.55) -- (2.3,-1.55);
    \draw[corr] (0.8,0.75) -- (1.9,1.55);
    \draw[corr] (0.8,0.75) -- (1.9,-1.55);
    \draw[corr] (3.0,-0.75) -- (1.9,1.55);
    \draw[corr] (3.0,-0.75) -- (1.9,-1.55);
    \draw[elec] (1.82,-1.75) -- (1.82,-1.35);
    \draw[elec] (1.98,-1.35) -- (1.98,-1.75);
    % d_z bracket between the two AO levels
    \draw[<->,cMuted,line width=0.7pt] (4.05,-0.75) -- (4.05,0.75);
    \node[cMuted,lab,anchor=west] at (4.12,0.0) {$d_z$};
    \draw[<->,cAccent,line width=0.7pt] (4.85,-1.55) -- (4.85,1.55);
    \node[cAccent,lab,anchor=west] at (4.95,0.0) {$|\dvec|$};
    \node[lab,anchor=north,align=center] at (1.9,-1.95)
      {$\sin^2\!\frac{\theta}{2}$ on A, $\cos^2\!\frac{\theta}{2}$ on B\\[-2pt]
       (polarized toward B)};
  \end{scope}
\end{tikzpicture}
\caption{Molecular orbital correlation diagrams for the two-orbital dimer (recall
$\beta<0$, so $\alpha+\beta$ is the \emph{lower} level). The splitting is
$|\dvec| = \sqrt{d_z^2+4\beta^2}$ in both cases; what changes is how the lower MO
distributes itself between the atoms, which is governed entirely by the mixing angle
$\theta$ of \cref{eq:mixing-angle}.}
\label{fig:mo-diagram}
\end{figure}
% ---------------------------------------------------------------------------

\begin{nbbox}[In the notebook]
Part 3 of the notebook builds this Hamiltonian, decomposes it, reproduces the numbers
above, and then scans $d_z$ to generate \cref{fig:avoided}. One of the test cells
asserts that the gap never falls below $2|\beta|$ across the entire scan --- a numerical
statement of the non-crossing rule, \cref{eq:noncrossing}.
\end{nbbox}

% ============================================================================
\section{Example II: a driven two-level system}
\label{sec:driven}
% ============================================================================

\subsection{The lab-frame Hamiltonian}

A spin-$1/2$ nucleus in a static field $B_0\hat{z}$ has the Zeeman Hamiltonian
%
\begin{equation}
\Hop_0 = -\gamma B_0 \Szo = \frac{\hbar\omega_0}{2}\szo,
\qquad \omega_0 \equiv -\gamma B_0 ,
\end{equation}
%
pure $\szo$, with the two spin states split by the Larmor energy $\hbar\omega_0$. For a
proton at \SI{11.7}{\tesla} this is \SI{500}{\mega\hertz}, the number printed on the
front of the spectrometer. Nothing happens under $\Hop_0$ alone: $\ket{\alpha}$ and
$\ket{\beta}$ are eigenstates and the Bloch vector merely spins about $\hat{z}$.

To make something happen, add a weak field rotating in the $xy$ plane at frequency
$\omega$ --- the RF coil in NMR, or, for an electronic transition in the dipole
approximation, the optical field of a laser:
%
\begin{equation}
\Hop(t) = \frac{\hbar\omega_0}{2}\szo
  + \frac{\hbar\omega_1}{2}\Big[\cos(\omega t+\varphi)\,\sxo + \sin(\omega t+\varphi)\,\syo\Big],
\qquad \omega_1 \equiv -\gamma B_1 .
\label{eq:lab-frame}
\end{equation}
%
This is genuinely time dependent, so the machinery of \cref{sec:geometry} does not apply
directly. One transformation fixes that.

\subsection{The rotating frame}

Move to a frame rotating about $z$ at the drive frequency by defining
$\ket{\tilde\psi} = \Ry(t)\ket{\psi}$ with $\Ry(t) = e^{+\ii\omega t\szo/2}$. Substituting
into the Schr\"odinger equation (worked out step by step in \cref{app:rotframe}) gives
an effective Hamiltonian $\tilde{H} = \Ry\Hop\Ry^\dagger - \tfrac{\hbar\omega}{2}\szo$
which turns out to be \emph{time independent}:

\begin{keyresult}[Rotating-frame Hamiltonian]
\vspace{-0.35cm}
\begin{equation}
\Hop_{\rm rot} = \frac{\hbar}{2}\Big[\underbrace{(\omega_0-\omega)}_{\textstyle \Delta}\,\szo
   + \omega_1\big(\cos\varphi\,\sxo + \sin\varphi\,\syo\big)\Big],
\label{eq:rotframe}
\end{equation}
that is, $d_z = \hbar\Delta$, $d_x = \hbar\omega_1\cos\varphi$,
$d_y = \hbar\omega_1\sin\varphi$.
\end{keyresult}

\noindent
Three features of \cref{eq:rotframe} are each a piece of real spectroscopic practice.

\begin{itemize}[leftmargin=*,itemsep=3pt]
\item \textbf{The detuning is the $\szo$ knob.} On resonance ($\omega=\omega_0$) it
vanishes and $\dvec$ lies flat in the equatorial plane. Off resonance, $\dvec$ tilts
toward the poles --- and by \cref{fig:bloch-prec}, a tilted $\dvec$ means the state can
no longer reach the far pole.

\item \textbf{The pulse phase is the $\syo$ knob.} Setting $\varphi=0$ gives a pure
$\sxo$ term, an ``$x$-pulse''; $\varphi=\pi/2$ gives a pure $\syo$ term, a ``$y$-pulse''.
This is not jargon: $\varphi$ is literally the phase of the RF pulse, and it is why every
pulse-sequence diagram in an NMR textbook labels its pulses $x$, $y$, $-x$, $-y$.

\item \textbf{The tilted-cone geometry of \cref{sec:bloch} is now physical.} The polar
angle of $\dvec$ is set by how far off resonance you are, and its azimuth by the pulse
phase.
\end{itemize}

\subsection{Rabi oscillations}

Apply \cref{eq:eigenvalues}. The splitting in the rotating frame defines the
\textbf{generalized Rabi frequency}
%
\begin{equation}
\Omega = \frac{|\dvec|}{\hbar} = \sqrt{\Delta^2+\omega_1^2},
\label{eq:rabi-freq}
\end{equation}
%
and integrating \cref{eq:bloch-eq} from $\svec(0)=\hat{z}$ (or equivalently applying
\cref{eq:eigenvectors} twice) gives

\begin{keyresult}[Rabi formula]
\vspace{-0.35cm}
\begin{equation}
P_\beta(t) = \frac{\omega_1^2}{\Omega^2}\,\sin^2\!\left(\frac{\Omega t}{2}\right).
\label{eq:rabi}
\end{equation}
\end{keyresult}

\noindent
\Cref{eq:rabi} repays careful reading. The \emph{rate} of oscillation,
$\Omega=\sqrt{\Delta^2+\omega_1^2}$, always \emph{increases} with detuning. But the
\emph{amplitude},
%
\begin{equation}
P_\beta^{\max} = \frac{\omega_1^2}{\Delta^2+\omega_1^2},
\label{eq:lorentzian}
\end{equation}
%
is a Lorentzian in $\Delta$ of half-width $\omega_1$: complete population inversion is
possible only exactly on resonance. \Cref{eq:lorentzian} is simultaneously the origin of
the spectroscopic lineshape and of \emph{pulse selectivity} --- a weak, long pulse
addresses a narrow band of frequencies; a strong, short pulse excites everything in
sight. \Cref{fig:rabi} shows both effects.

% ---------------------------------------------------------------------------
\begin{figure}[t]
\centering
\begin{tikzpicture}
\begin{groupplot}[
  group style={group size=2 by 1, horizontal sep=1.5cm},
  notestyle, height=6.0cm,
]
\nextgroupplot[
  xlabel={$\omega_1 t$}, ylabel={$P_\beta(t)$},
  title={Rabi oscillations},
  xmin=0, xmax=12.566, ymin=-0.03, ymax=1.05,
  xtick={0,3.1416,6.2832,9.4248,12.566},
  xticklabels={$0$,$\pi$,$2\pi$,$3\pi$,$4\pi$},
  legend style={at={(0.5,-0.30)},anchor=north,legend columns=4,draw=none,fill=none,
                font=\footnotesize},
  legend cell align=left,
]
\addplot[cLower, very thick, domain=0:12.566, samples=400]
  {(1/1)*sin(deg(sqrt(1)*x/2))^2};                       \addlegendentry{$\Delta=0$}
\addplot[cAccent, very thick, domain=0:12.566, samples=400]
  {(1/2)*sin(deg(sqrt(2)*x/2))^2};                       \addlegendentry{$\Delta=\omega_1$}
\addplot[cWarm, very thick, domain=0:12.566, samples=500]
  {(1/5)*sin(deg(sqrt(5)*x/2))^2};                       \addlegendentry{$\Delta=2\omega_1$}
\addplot[cUpper, very thick, domain=0:12.566, samples=600]
  {(1/17)*sin(deg(sqrt(17)*x/2))^2};                     \addlegendentry{$\Delta=4\omega_1$}

\nextgroupplot[
  xlabel={detuning $\Delta/\omega_1$}, ylabel={$P_\beta^{\max}$},
  title={\ldots but only on resonance is inversion complete},
  xmin=-5, xmax=5, ymin=-0.03, ymax=1.06,
]
\addplot[cUpper, very thick, domain=-5:5, samples=300]{1/(1+x^2)};
\addplot[cMuted, dotted, thin, domain=-5:5, samples=2]{0.5};
\draw[<->,cAccent,thin] (axis cs:-1,0.5) -- (axis cs:1,0.5);
\node[cAccent,font=\footnotesize,anchor=south] at (axis cs:0,0.53){FWHM $=2\omega_1$};
\end{groupplot}
\end{tikzpicture}
\caption{\textbf{Left:} \cref{eq:rabi} at four detunings. Detuning makes the
oscillation \emph{faster} and \emph{shallower} at the same time. \textbf{Right:} the
maximum transferable population, \cref{eq:lorentzian}. A pulse of Rabi strength
$\omega_1$ can only fully invert spins within roughly $\pm\omega_1$ of resonance ---
which is what ``selective pulse'' means.}
\label{fig:rabi}
\end{figure}
% ---------------------------------------------------------------------------

\subsection{$\pi$ and $\pi/2$ pulses}

On resonance $\dvec$ lies in the equatorial plane, the precession cone opens into a
great circle through both poles, and the state sweeps from $\ket{\alpha}$ to
$\ket{\beta}$ and back. Two durations matter:
%
\begin{equation}
t_\pi = \frac{\pi}{\omega_1}: \;\; \ket{\alpha}\to\ket{\beta} \text{ (inversion)},
\qquad
t_{\pi/2} = \frac{\pi}{2\omega_1}: \;\; \ket{\alpha}\to \text{the equator (coherence)}.
\end{equation}
%
The $\pi/2$ pulse is the one that generates an NMR signal: it converts population
difference into transverse coherence, which is what the detection coil sees as the free
induction decay. And the pulse phase decides \emph{where} on the equator you land. Using
\cref{eq:bloch-eq} with $\dvec = \hbar\omega_1\hat{x}$,
%
\begin{equation}
\text{$x$-pulse: } \hat{z}\to-\hat{y},
\qquad
\text{$y$-pulse: } \hat{z}\to+\hat{x},
\end{equation}
%
two states $\ang{90}$ apart on the equator. Both are ``the same pulse'' in every respect
except phase.

\begin{lookahead}[Where this is going: the Jaynes--Cummings model]
Write the transverse part of \cref{eq:rotframe} using the ladder operators
$\hat\sigma_\pm = \half(\sxo\pm\ii\syo)$:
\begin{equation}
\frac{\hbar\omega_1}{2}\big(\cos\varphi\,\sxo+\sin\varphi\,\syo\big)
 = \frac{\hbar\omega_1}{2}\Big(e^{-\ii\varphi}\hat\sigma_+ + e^{+\ii\varphi}\hat\sigma_-\Big).
\end{equation}
The drive \emph{raises} the two-level system while supplying a phase, or \emph{lowers} it
while absorbing one. Here the field is a classical number: $\omega_1\propto B_1$, and it
is inexhaustible --- the system can cycle forever and the field never notices.

Now promote the field to a quantum degree of freedom. The phase factors become photon
annihilation and creation operators, $e^{-\ii\varphi}\to\hat{a}$ and
$e^{+\ii\varphi}\to\hat{a}^\dagger$, and
\begin{equation}
\Hop_{\rm JC} = \frac{\hbar\omega_0}{2}\szo + \hbar\omega\,\hat{a}^\dagger\hat{a}
  + \hbar g\big(\hat{a}\,\hat\sigma_+ + \hat{a}^\dagger\hat\sigma_-\big),
\label{eq:jc}
\end{equation}
the \textbf{Jaynes--Cummings} Hamiltonian \cite{JaynesCummings} --- the single-emitter
case of the Tavis--Cummings model built in Computational Set 4. The Hilbert space is no
longer two dimensional; it is spin $\otimes$ boson. The classical drive strength
$\omega_1$ is replaced by $2g\sqrt{n+1}$, so that even the vacuum ($n=0$) drives Rabi
oscillations. That is vacuum Rabi splitting, and it is the foundation of polaritonic
chemistry. Everything in this section is its classical-field limit.
\end{lookahead}

\begin{nbbox}[In the notebook]
Part 5 verifies \cref{eq:rabi} against numerical propagation to $10^{-9}$ for several
detunings, confirms that a resonant $\pi$ pulse inverts the population exactly, and
shows that the $x$- and $y$-pulses leave the Bloch vector $\ang{90}$ apart.
\end{nbbox}

% ============================================================================
\section{Example III: donor--acceptor electron transfer}
\label{sec:et}
% ============================================================================

\subsection{Diabatic states and a nuclear coordinate}

Consider an electron transferring from a donor to an acceptor,
$\Dm \longrightarrow \Am$. The two basis states are the \textbf{diabatic}
(charge-localized) states $\ket{D}\equiv\ket{\alpha}$ and $\ket{A}\equiv\ket{\beta}$.

What is new here is that their energies are not fixed numbers. Solvent molecules
reorganize around whichever site currently holds the charge, so the diabatic energies
depend on where the nuclei are. Introduce a single collective nuclear coordinate $q$ ---
the solvent or reaction coordinate --- and model each diabatic surface as a harmonic
well of the same force constant $k$, displaced by $\pm q_0$:
%
\begin{equation}
V_D(q) = \half k(q+q_0)^2,
\qquad
V_A(q) = \half k(q-q_0)^2 + \Delta G^{\circ}.
\label{eq:diabats}
\end{equation}
%
The donor state is happiest at $q=-q_0$, the acceptor state at $q=+q_0$, and
$\Delta G^\circ$ is the thermodynamic driving force. The electronic coupling
$V_{DA} = \braket{D|\Hop|A}$ between them is essentially independent of $q$.

\subsection{Identifying $\dvec$}

With $\Hop(q)$ built from \cref{eq:diabats},
%
\begin{equation}
d_z(q) = V_D(q)-V_A(q) = 2kq_0\,q - \Delta G^\circ,
\qquad
d_x = 2V_{DA},
\qquad d_y = 0.
\end{equation}
%
Defining the \textbf{reorganization energy} $\lambda \equiv 2kq_0^2$ --- the energy cost
of distorting the nuclei from the donor's optimal geometry to the acceptor's
\emph{without} moving the electron:

\begin{keyresult}[Electron transfer]
\vspace{-0.35cm}
\begin{equation}
d_z(q) = \frac{\lambda}{q_0}\,q - \Delta G^\circ,
\qquad
d_x = 2V_{DA}.
\label{eq:et-d}
\end{equation}
\end{keyresult}

\noindent
This is the structural heart of the example, and it is worth stating plainly:
\textbf{the $\szo$ coefficient is linear in a nuclear coordinate, and the $\sxo$
coefficient is a constant.} Compare \cref{eq:dimer-d}, where $d_z$ was a fixed
electronegativity difference. Here the molecule sweeps its own $d_z$ back and forth as
it vibrates, so the parameter scan of \cref{fig:avoided} is not a hypothetical --- it is
something the molecule physically does, many times per picosecond.

\subsection{Adiabatic surfaces and the Marcus barrier}

The eigenvalues of $\Hop(q)$ are the \textbf{adiabatic} potential energy surfaces --- the
Born--Oppenheimer curves the nuclei actually move on:
%
\begin{equation}
E_\pm(q) = \bar{V}(q) \pm \half\sqrt{d_z(q)^2 + 4V_{DA}^2},
\qquad \bar{V}(q) = \half\!\left[V_D(q)+V_A(q)\right].
\label{eq:adiabats}
\end{equation}
%
Three consequences follow at once from \cref{sec:geometry}:

\begin{itemize}[leftmargin=*,itemsep=2pt]
\item the diabats cross where $d_z(q^\ddagger)=0$, that is at
      $q^\ddagger = \Delta G^\circ q_0/\lambda$;
\item the adiabatic gap there is exactly $|d_x| = 2V_{DA}$;
\item the activation energy, measured on the \emph{diabatic} surface from the donor
      minimum up to the crossing, is
      \begin{equation}
      \Delta G^{\ddagger} = V_D(q^\ddagger) - V_D(-q_0)
          = \frac{(\lambda+\Delta G^\circ)^2}{4\lambda},
      \label{eq:marcus-barrier}
      \end{equation}
      the \textbf{Marcus} activation free energy \cite{Marcus1956,MarcusNobel}
      (algebra in \cref{app:marcus}).
\end{itemize}

\Cref{eq:marcus-barrier} is a parabola in the driving force, minimized --- and equal to
zero --- at $\Delta G^\circ = -\lambda$. Since the nonadiabatic rate goes as
$k_{ET}\propto |V_{DA}|^2\exp(-\Delta G^\ddagger/k_BT)$, making a reaction \emph{more}
exergonic speeds it up only until $-\Delta G^\circ = \lambda$, and slows it down
thereafter. This is the celebrated \textbf{inverted region}, whose experimental
confirmation was central to the 1992 Nobel Prize in Chemistry. \Cref{fig:et,fig:marcus}
show the surfaces and the rate.

% ---------------------------------------------------------------------------
\begin{figure}[t]
\centering
\begin{tikzpicture}
\begin{axis}[
  widenote, height=7.4cm,
  xlabel={nuclear (solvent) coordinate $q$},
  ylabel={energy},
  xmin=-2.25, xmax=2.25, ymin=-0.95, ymax=1.75,
  legend style={at={(0.5,-0.22)},anchor=north,legend columns=2,draw=none,fill=none,
                font=\footnotesize},
  legend cell align=left,
]
% adiabats: soft thick bands underneath
\addplot[cLower, line width=3.6pt, opacity=0.35, domain=-2.25:2.25, samples=400]
  {0.5*(0.5*(x+1)^2 + 0.5*(x-1)^2 - 0.5) - 0.5*sqrt((2*x+0.5)^2 + 0.09)};
\addlegendentry{adiabat $E_-$}
\addplot[cUpper, line width=3.6pt, opacity=0.35, domain=-2.25:2.25, samples=400]
  {0.5*(0.5*(x+1)^2 + 0.5*(x-1)^2 - 0.5) + 0.5*sqrt((2*x+0.5)^2 + 0.09)};
\addlegendentry{adiabat $E_+$}
% diabats on top, crisp
\addplot[cMuted, dashed, thick, domain=-2.25:2.25, samples=200]{0.5*(x+1)^2};
\addlegendentry{diabat $V_D$ (charge on donor)}
\addplot[cMuted, densely dotted, thick, domain=-2.25:2.25, samples=200]{0.5*(x-1)^2-0.5};
\addlegendentry{diabat $V_A$ (charge on acceptor)}
% avoided crossing marker at q* = -0.25
\draw[<->,cAccent,thick] (axis cs:-0.25,0.13125) -- (axis cs:-0.25,0.43125);
\node[cAccent,font=\footnotesize,anchor=west] at (axis cs:-0.17,0.60){$2V_{DA}$};
% Marcus barrier
\draw[<->,cWarm,thick] (axis cs:-1.28,0.0) -- (axis cs:-1.28,0.2812);
\draw[cWarm,densely dash dot,thin] (axis cs:-1.36,0.2812) -- (axis cs:-0.25,0.2812);
\draw[cWarm,densely dash dot,thin] (axis cs:-1.36,0.0) -- (axis cs:-1.0,0.0);
\node[cWarm,font=\footnotesize,anchor=east,fill=white,inner sep=1pt]
  at (axis cs:-1.34,0.145){$\Delta G^{\ddagger}$};
% minima
\addplot[only marks, mark=*, mark size=1.6pt, black] coordinates {(-1,0) (1,-0.5)};
\node[font=\footnotesize,anchor=north] at (axis cs:-1,-0.07){$\Dm$};
\node[font=\footnotesize,anchor=north] at (axis cs:1,-0.57){$\Am$};
% driving force
\draw[<->,cPurple,thin] (axis cs:1.62,0.0) -- (axis cs:1.62,-0.5);
\node[cPurple,font=\footnotesize,anchor=west] at (axis cs:1.66,-0.25){$\Delta G^{\circ}$};
\draw[cPurple,densely dotted,thin] (axis cs:-1,0) -- (axis cs:1.62,0);
\end{axis}
\end{tikzpicture}
\caption{Diabatic and adiabatic surfaces for
$k=1$, $q_0=1$ (so $\lambda=2$), $\Delta G^\circ = -0.5$, $V_{DA}=0.15$. Away from the
crossing the adiabats and diabats coincide --- the electron stays put. Near
$q^\ddagger = \Delta G^\circ q_0/\lambda = -0.25$ the coupling splits them by $2V_{DA}$
and the character of each adiabatic surface switches over.}
\label{fig:et}
\end{figure}
% ---------------------------------------------------------------------------

% ---------------------------------------------------------------------------
\begin{figure}[t]
\centering
\begin{tikzpicture}
\begin{groupplot}[
  group style={group size=2 by 1, horizontal sep=1.5cm},
  notestyle, height=5.8cm,
]
\nextgroupplot[
  xlabel={$\Delta G^\circ/\lambda$}, ylabel={$\Delta G^{\ddagger}/\lambda$},
  title={Activation barrier},
  xmin=-3, xmax=0.5, ymin=-0.03, ymax=1.05,
]
\addplot[cWarm, very thick, domain=-3:0.5, samples=200]{(1+x)^2/4};
\draw[cMuted,densely dotted,thin] (axis cs:-1,-0.03) -- (axis cs:-1,1.05);
\node[cMuted,font=\footnotesize,align=center] at (axis cs:-2.1,0.62){inverted\\region};
\node[cMuted,font=\footnotesize,align=center] at (axis cs:-0.35,0.62){normal\\region};

\nextgroupplot[
  xlabel={$\Delta G^\circ/\lambda$}, ylabel={rate $\propto e^{-\Delta G^{\ddagger}/k_BT}$},
  title={Marcus rate ($k_BT=\lambda/20$)},
  xmin=-3, xmax=0.5, ymode=log, ymin=1e-9, ymax=4,
  ytick={1e-8,1e-6,1e-4,1e-2,1},
]
\addplot[cPurple, very thick, domain=-3:0.5, samples=300]{exp(-5*(1+x)^2)};
\draw[cMuted,densely dotted,thin] (axis cs:-1,1e-9) -- (axis cs:-1,4);
\end{groupplot}
\end{tikzpicture}
\caption{The Marcus inverted region. The barrier \eqref{eq:marcus-barrier} is a parabola
in the driving force, so the rate is \emph{not} monotonic: past
$-\Delta G^\circ=\lambda$ a more exergonic reaction is a slower one.}
\label{fig:marcus}
\end{figure}
% ---------------------------------------------------------------------------

\subsection{Landau--Zener: sweeping through the crossing}

Everything above is static --- surfaces evaluated at fixed $q$. Now let the nuclei
actually move through the crossing. If the coordinate sweeps linearly in time, $q = vt$,
then by \cref{eq:et-d}
%
\begin{equation}
d_z(t) = \alpha t, \qquad \alpha = \frac{\lambda v}{q_0}
  = \left|\frac{\dd}{\dd t}\big(V_D-V_A\big)\right| ,
\end{equation}
%
the rate at which the diabatic gap is swept. Starting far to one side in a single
diabatic state, two outcomes compete:

\begin{itemize}[leftmargin=*,itemsep=2pt]
\item \textbf{adiabatic passage} --- the system follows the lower adiabatic surface and
therefore \emph{changes} diabatic character: the electron transfers;
\item \textbf{diabatic passage} --- the system blasts straight through the avoided
crossing, staying on the same diabat: no transfer.
\end{itemize}

\noindent
Landau and Zener solved this exactly in 1932 \cite{Zener1932,Landau1932}:

\begin{keyresult}[Landau--Zener]
\vspace{-0.35cm}
\begin{equation}
P_{\rm diabatic} = \exp\!\left(-\frac{2\pi V_{DA}^2}{\hbar\,\alpha}\right).
\label{eq:lz}
\end{equation}
\end{keyresult}

\noindent
Slow sweeps or strong coupling give adiabatic passage ($P\to0$); fast sweeps or weak
coupling give diabatic passage ($P\to1$). Note the combination in the exponent:
$(\text{coupling})^2$ divided by (sweep rate). Expanding for small $V_{DA}$,
%
\begin{equation}
1 - P_{\rm diabatic} \approx \frac{2\pi V_{DA}^2}{\hbar\alpha},
\end{equation}
%
which is the same $|V_{DA}|^2$ proportionality that appears in Fermi's golden rule and
in the nonadiabatic Marcus rate --- the two results are the same physics seen from two
directions.

% ---------------------------------------------------------------------------
\begin{figure}[t]
\centering
\begin{tikzpicture}
\begin{axis}[
  notestyle, width=0.62\textwidth, height=6.2cm,
  xlabel={diabatic coupling $V_{DA}$},
  ylabel={$P_{\rm diabatic}$},
  xmin=0, xmax=0.53, ymin=0.13, ymax=1.05,
  legend pos=north east, legend cell align=left,
]
\addplot[black, very thick, domain=0:0.53, samples=200]{exp(-2*pi*x^2)};
\addlegendentry{Zener formula, Eq.~(\ref{eq:lz})}
\addplot[only marks, mark=*, mark size=2.1pt, cUpper]
  table[x index=0, y index=1] {figures/landau_zener_numeric.dat};
\addlegendentry{numerical propagation}
\end{axis}
\end{tikzpicture}
\caption{\Cref{eq:lz} against direct numerical propagation through the crossing
($\hbar=1$, sweep rate $\alpha=1$). The numerical points come from the notebook, which
slices the sweep into 60000 steps and evaluates $\Hop$ at the midpoint of each. The
residual disagreement, at most $0.004$, is not step error: Zener's result is exact only
for a sweep running from $t=-\infty$ to $+\infty$, and any truncated sweep leaves a
little interference behind.}
\label{fig:lz}
\end{figure}
% ---------------------------------------------------------------------------

\begin{lookahead}[Where this is going: the Holstein model]
Collecting \cref{eq:et-d},
\begin{equation}
\Hop(q) = \bar{V}(q)\,\Iop
 + \half\!\left[\frac{\lambda}{q_0}q - \Delta G^\circ\right]\szo + V_{DA}\,\sxo .
\end{equation}
We treated $q$ as a \emph{classical parameter} --- a number to dial, or to sweep linearly
in time. But $q$ is a nuclear coordinate, and nuclei are quantum. Promote it to a
harmonic oscillator, $\hat{q}\propto(\hat{a}+\hat{a}^\dagger)$, and
\begin{equation}
\Hop_{\rm Holstein} = \frac{\varepsilon}{2}\szo + V_{DA}\,\sxo
  + \hbar\omega\,\hat{a}^\dagger\hat{a}
  + \hbar\omega g\,(\hat{a}+\hat{a}^\dagger)\,\szo ,
\label{eq:holstein}
\end{equation}
the \textbf{Holstein}, or spin--boson, Hamiltonian \cite{Holstein1959,Leggett1987}.

Set \cref{eq:holstein} beside \cref{eq:jc} and notice how different they are despite
sharing a two-dimensional electronic space and one boson:
\begin{center}
\small
\begin{tabular}{@{}lll@{}}
\toprule
Model & Coupling operator & The boson\ldots \\
\midrule
Jaynes--Cummings & $\hat{a}\hat\sigma_+ + \hat{a}^\dagger\hat\sigma_-$
  & \textbf{flips} the two-level system \\
Holstein & $(\hat{a}+\hat{a}^\dagger)\,\szo$
  & \textbf{shifts} its energies \\
\bottomrule
\end{tabular}
\end{center}
One gives polaritons and vacuum Rabi splitting; the other gives vibronic progressions,
Franck--Condon factors, and polaron formation. You have now built the single-spin
ancestor of both.
\end{lookahead}

% ============================================================================
\section{Example IV: the ammonia inversion doublet}
\label{sec:ammonia}
% ============================================================================

The three examples so far built their two-level space from two places an electron can
be. This one builds it from two \emph{shapes the same molecule can have}.

\subsection{Tunneling as an off-diagonal matrix element}

Ammonia is pyramidal, and the nitrogen can pass through the plane of the three hydrogens
like an umbrella turning inside out. Label the two pyramidal geometries $\ket{L}$
(umbrella down) and $\ket{R}$ (umbrella up). Classically they are separated by a barrier
of roughly \SI{2000}{\wavenumber} and the molecule would be stuck in one of them.
Quantum mechanically the nitrogen \textbf{tunnels}, and tunneling is nothing other than
an off-diagonal matrix element:
%
\begin{equation}
\Hop = \begin{bmatrix} 0 & -\Delta_t \\ -\Delta_t & 0\end{bmatrix} = -\Delta_t\,\sxo .
\label{eq:nh3-bare}
\end{equation}
%
Pure $\sxo$. The eigenstates are $\tfrac{1}{\sqrt2}(\ket{L}\pm\ket{R})$, split by
$2\Delta_t$: the \textbf{inversion doublet}, measured at
%
\begin{equation}
2\Delta_t = \SI{0.7934}{\wavenumber} = \SI{23.786}{\giga\hertz},
\end{equation}
%
a microwave transition --- the line Gordon, Zeiger and Townes used to build the first
\textbf{maser} in 1954, three years before the laser \cite{Townes1955}.

The conceptual point is worth dwelling on. The energy eigenstates of ammonia are
\emph{not} the pyramidal geometries; they are delocalized over both. A molecule prepared
as $\ket{L}$ is a superposition of the two eigenstates, so by \cref{sec:bloch} it
oscillates: the umbrella flips back and forth at \SI{23.786}{\giga\hertz}.

\subsection{Adding a $\szo$ term: the Stark effect}

Switch on a static electric field $\mathcal{E}$ along the molecular axis. The two
pyramidal geometries have dipole moments pointing in \emph{opposite} directions, so the
field stabilizes one and destabilizes the other:

\begin{keyresult}[Ammonia in a field]
\vspace{-0.35cm}
\begin{equation}
\Hop = \mu\mathcal{E}\,\szo - \Delta_t\,\sxo
\qquad\Longrightarrow\qquad
d_z = 2\mu\mathcal{E}, \quad d_x = -2\Delta_t ,
\label{eq:nh3-stark}
\end{equation}
with $\mu = \SI{1.47}{\debye}$ the NH$_3$ dipole moment.
\end{keyresult}

\noindent
Structurally this is \cref{eq:dimer-d} again --- but now \emph{you} control $d_z$ with a
knob on the bench. The gap is
%
\begin{equation}
E_+-E_- = \sqrt{(2\mu\mathcal{E})^2+(2\Delta_t)^2},
\end{equation}
%
quadratic in $\mathcal{E}$ at low field (a second-order Stark effect) and crossing over
to linear at high field. The crossover occurs where the two terms are comparable,
$\mu\mathcal{E}=\Delta_t$, that is at
%
\begin{equation}
\mathcal{E}_{\rm cross} = \frac{\Delta_t}{\mu} = \SI{16.07}{\kilo\volt\per\centi\meter},
\end{equation}
%
at which point the gap has grown by exactly $\sqrt{2}$. Turn the field well past this
and $\theta\to0$: the field literally traps the umbrella on one side, converting a
delocalized molecule into a localized one. \Cref{fig:nh3} shows both halves.

% ---------------------------------------------------------------------------
\begin{figure}[t]
\centering
\begin{tikzpicture}
\begin{groupplot}[
  group style={group size=2 by 1, horizontal sep=1.6cm},
  notestyle, height=6.0cm,
  xmin=0, xmax=60,
]
\nextgroupplot[
  xlabel={electric field $\mathcal{E}$ (kV\,cm$^{-1}$)},
  ylabel={energy (cm$^{-1}$)},
  title={Stark effect on the doublet},
  ymin=-1.7, ymax=1.7,
  legend style={at={(0.03,0.97)},anchor=north west,draw=none,fill=white,
                fill opacity=0.85,text opacity=1,font=\footnotesize},
  legend cell align=left,
]
% mu = 0.024685 cm^-1 per kV/cm ; Delta_t = 0.3967 cm^-1
\addplot[cUpper, very thick, domain=0:60, samples=200]
  {sqrt((0.024685*x)^2 + 0.3967^2)};
\addlegendentry{upper level}
\addplot[cLower, very thick, domain=0:60, samples=200]
  {-sqrt((0.024685*x)^2 + 0.3967^2)};
\addlegendentry{lower level}
\addplot[cMuted, dashed, thin, domain=0:60, samples=2]{0.024685*x};
\addlegendentry{$\pm\mu\mathcal{E}$ (no tunneling)}
\addplot[cMuted, dashed, thin, domain=0:60, samples=2]{-0.024685*x};
\draw[cMuted,densely dotted,thin] (axis cs:16.07,-1.7) -- (axis cs:16.07,1.7);
\draw[<->,cAccent,thick] (axis cs:2.6,-0.3967) -- (axis cs:2.6,0.3967);
\node[cAccent,font=\footnotesize,anchor=west] at (axis cs:3.2,0.72){$2\Delta_t$};

\nextgroupplot[
  xlabel={electric field $\mathcal{E}$ (kV\,cm$^{-1}$)},
  ylabel={$|\Braket{\szo}|$ of the lower state},
  title={The field traps the umbrella},
  ymin=-0.03, ymax=1.03,
]
\addplot[cAccent, very thick, domain=0:60, samples=300]
  {abs(0.024685*x)/sqrt((0.024685*x)^2 + 0.3967^2)};
\draw[cMuted,densely dotted,thin] (axis cs:16.07,-0.03) -- (axis cs:16.07,1.03);
\node[cMuted,font=\footnotesize,anchor=west] at (axis cs:17.5,0.20)
  {$\mu\mathcal{E}=\Delta_t$};
\node[font=\footnotesize,anchor=west,align=left] at (axis cs:1.5,0.86)
  {delocalized $\rightarrow$ localized};
\end{groupplot}
\end{tikzpicture}
\caption{Ammonia in a static electric field, from \cref{eq:nh3-stark}. \textbf{Left:}
the levels start split by the tunneling doublet $2\Delta_t = \SI{0.7934}{\wavenumber}$
and approach the classical Stark asymptotes $\pm\mu\mathcal{E}$ at high field.
\textbf{Right:} the localization $|\Braket{\szo}| = |d_z|/|\dvec| = |\cos\theta|$ of the
lower state, running from $0$ (equal amplitude on both umbrella geometries) to $1$
(trapped on one side).}
\label{fig:nh3}
\end{figure}
% ---------------------------------------------------------------------------

\begin{nbbox}[A puzzle worth posing to yourself]
At zero field the ground state is $\tfrac{1}{\sqrt2}(\ket{L}+\ket{R})$, which has
\emph{zero} dipole moment --- yet NH$_3$ is a famously polar molecule with
$\mu=\SI{1.47}{\debye}$. Both statements are true. The resolution turns on what a
dipole-moment measurement actually averages over, and on what timescale; see
Problem~\ref{prob:nh3} in \cref{sec:problems}.
\end{nbbox}

% ============================================================================
\section{Where does $\syo$ come from?}
\label{sec:sigmay}
% ============================================================================

Of the four examples, only the driven system of \cref{sec:driven} carried a $\syo$ term.
The LCAO dimer, the electron-transfer complex and ammonia all had $d_y=0$. That is not a
coincidence, and understanding why is worth a short section.

\subsection{For one static Hamiltonian, $\syo$ is a choice of convention}

The overall phase of a basis vector is a convention, not physics. Rephase the second
basis state,
%
\begin{equation}
\ket{\beta} \;\longrightarrow\; e^{\ii\chi}\ket{\beta},
\qquad\text{i.e.}\qquad
\hat{U} = \mathrm{diag}\!\left(1, e^{\ii\chi}\right),
\qquad
\Hop' = \hat{U}^\dagger \Hop\, \hat{U}.
\label{eq:rephase}
\end{equation}
%
This leaves the diagonal untouched and multiplies the off-diagonal element by
$e^{\ii\chi}$. Since $H_{\alpha\beta} = \half(d_x-\ii d_y) = \half|\dperp|e^{-\ii\varphi}$,
choosing $\chi=\varphi$ makes $H'_{\alpha\beta}$ real and positive:
%
\begin{equation}
d_x' = \dperp = \sqrt{d_x^2+d_y^2}, \qquad d_y' = 0, \qquad d_z' = d_z .
\end{equation}
%
The transverse components have simply been rotated into one another. The spectrum,
being unitarily invariant, is unchanged --- as it must be, since \cref{eq:eigenvalues}
depends on $d_x$ and $d_y$ only through $\dperp$.

\begin{keyresult}[The gauge statement]
For a \emph{single, static} two-level Hamiltonian, a $\syo$ term carries no physical
content: it can always be rephased away, leaving a real symmetric matrix. Every
observable depends on $d_x$ and $d_y$ only through $\dperp = \sqrt{d_x^2+d_y^2}$.
\end{keyresult}

\subsection{Three ways to make it physical}

$\syo$ acquires meaning only \emph{relative to something else} that fixes the phase
convention. There are three such situations, and each is worth recognizing on sight.

\paragraph{1. A time-dependent drive.} In \cref{sec:driven} the pulse phase $\varphi$
moved $\dvec$ around the equatorial plane. Any \emph{single} pulse may be declared an
``$x$-pulse'' by fiat --- but once that choice is made, the phase of the \emph{second}
pulse in a sequence is physical and measurable. This is exactly what two-dimensional NMR
and Ramsey interferometry exploit (Problem~\ref{prob:ramsey}).

\paragraph{2. Broken time-reversal symmetry.} For a spinless electronic Hamiltonian in a
basis of real orbitals with no magnetic field, time-reversal symmetry forces $\Hop$ to be
real symmetric, so $d_y=0$ automatically --- which is why \cref{sec:dimer,sec:et} had
none. To get an unremovable $\syo$ you must break that symmetry: a magnetic field, which
attaches a Peierls phase $\beta\to\beta e^{\ii\phi_B}$ to each hopping integral;
circularly polarized light; or \textbf{spin--orbit coupling}, which is why the effective
$2\times2$ Hamiltonian of a Kramers doublet in a heavy-atom complex generally carries all
three Pauli components.

\paragraph{3. Three or more coupled states.} This is the subtle one. With $N$ basis states
you have $N-1$ useful rephasings available (the overall phase does nothing), so you can
kill at most $N-1$ hopping phases. A ring of three sites has three bonds and only two
usable phases: one combination survives. What survives is the phase accumulated around
the closed loop --- a magnetic flux --- and it is genuinely gauge invariant. This is the
origin of the Aharonov--Bohm effect and, in a different guise, of the Berry phase.
Problem~\ref{prob:peierls} works out the dimer and the triangle side by side.

\begin{nbbox}[The practical rule]
When you meet a real electronic-structure Hamiltonian in a real orbital basis, expect
$\szo$ and $\sxo$ only. A $\syo$ term is a signal that a field, a spin--orbit
interaction, a time-dependent drive, or a closed loop is somewhere in the problem.
\end{nbbox}

% ============================================================================
\section{Summary}
\label{sec:summary}
% ============================================================================

Every system in these notes is the same equation,
$\Hop = \varepsilon_0\Iop + \half\dvec\cdot\pauliv$. Only the dictionary changes.

\begin{table}[h]
\centering
\small
\setlength{\tabcolsep}{5pt}
\begin{tabular}{@{}p{2.55cm}p{2.75cm}p{3.1cm}p{2.5cm}p{2.9cm}@{}}
\toprule
\textbf{System} & $\ket{\alpha},\ket{\beta}$ & $d_z$ (asymmetry) & $d_x,d_y$ (coupling)
  & $|\dvec|$ is called \\
\midrule
LCAO dimer\newline(\cref{sec:dimer})
  & AO on atom A / atom B
  & $\alpha_A-\alpha_B$\newline (electronegativity)
  & $2\beta$\newline (resonance integral)
  & bonding--antibonding gap \\[4pt]
Driven spin\newline(\cref{sec:driven})
  & spin up / down
  & $\hbar(\omega_0-\omega)$\newline (detuning)
  & $\hbar\omega_1 e^{\pm\ii\varphi}$\newline (drive $\times$ phase)
  & $\hbar\Omega$, generalized Rabi \\[4pt]
Electron transfer\newline(\cref{sec:et})
  & charge on D / on A
  & $(\lambda/q_0)q-\Delta G^\circ$\newline \emph{(nuclear coordinate!)}
  & $2V_{DA}$\newline (electronic coupling)
  & adiabatic gap \\[4pt]
NH$_3$ inversion\newline(\cref{sec:ammonia})
  & umbrella down / up
  & $2\mu\mathcal{E}$\newline (Stark field)
  & $-2\Delta_t$\newline (tunneling)
  & inversion doublet \\
\bottomrule
\end{tabular}
\caption{One Hamiltonian, four molecules.}
\label{tab:dictionary}
\end{table}

\noindent
And the physics is always the same three statements:

\begin{enumerate}[leftmargin=*,itemsep=3pt]
\item \textbf{Splitting.} $E_\pm = \varepsilon_0 \pm \half|\dvec|$. Because $d_z$ and
$\dperp$ add in quadrature, coupled levels never cross; the closest they come is
$\dperp$.
\item \textbf{Mixing.} $\tan\theta = \dperp/d_z$ decides whether the eigenstates are
localized on the basis states or delocalized between them. The switch happens over a
window of width $\sim\dperp$.
\item \textbf{Dynamics.} $\dot{\svec} = \hbar^{-1}\dvec\times\svec$. The state precesses
about $\dvec$ at the Bohr frequency $|\dvec|/\hbar$.
\end{enumerate}

\noindent
Faced with a new two-level problem --- an exciton dimer, a superconducting qubit, a
Kramers doublet, a curve crossing in a photochemical reaction --- the first move is
always to find $\dvec$.

\subsection*{What comes next}

Both models flagged along the way live in a \emph{larger} Hilbert space: spin
$\otimes$ boson, not spin alone. That tensor product is the step these notes deliberately
do not take.

\begin{center}
\small
\begin{tabular}{@{}lll@{}}
\toprule
Model & Coupling term & Its single-spin ancestor \\
\midrule
Jaynes--Cummings, \cref{eq:jc}
  & $\hbar g(\hat{a}\hat\sigma_+ + \hat{a}^\dagger\hat\sigma_-)$
  & \cref{sec:driven}, field as a classical number \\
Holstein, \cref{eq:holstein}
  & $\hbar\omega g(\hat{a}+\hat{a}^\dagger)\szo$
  & \cref{sec:et}, $q$ as a classical parameter \\
\bottomrule
\end{tabular}
\end{center}

\noindent
In both cases the two-level machinery developed here survives intact --- it is simply
tensored with a harmonic oscillator. Keep the results boxed in
\cref{sec:formalism,sec:geometry}; you will need every one of them again.

% ============================================================================
\section{Problems}
\label{sec:problems}
% ============================================================================

These are the problems that close the notebook.\ifsolutions{} Numerical answers quoted
in the solutions have been checked against direct diagonalization or propagation.\fi

% ---------------------------------------------------------------------------
\begin{problem}{The excitonic dimer}
\label{prob:exciton}
Two chromophores sit close enough that their transition dipoles interact. In the
singly-excited basis $\{\ket{e_1 g_2},\ket{g_1 e_2}\}$ the Hamiltonian is
\[
\Hop = \begin{bmatrix}\varepsilon_1 & J \\ J & \varepsilon_2\end{bmatrix},
\]
with $J$ the excitonic coupling.
\begin{enumerate}[label=(\alph*),leftmargin=*,itemsep=1pt]
\item Identify $\varepsilon_0$, $d_z$ and $d_x$.
\item For identical chromophores, find the eigenstates and energies. The splitting
$2|J|$ is the \emph{Davydov splitting}.
\item The transition dipole to the symmetric state is
$\bm{\mu}_1+\bm{\mu}_2$ and to the antisymmetric state $\bm{\mu}_1-\bm{\mu}_2$. In the
point-dipole approximation, $J>0$ for a side-by-side (H-aggregate) geometry and $J<0$
for head-to-tail (J-aggregate). In each case, is the \emph{bright} state the upper or
lower one, and does the absorption red- or blue-shift relative to the monomer?
\item Adapt the code from Part 3 of the notebook to plot the two exciton energies and
their oscillator strengths against $\varepsilon_1-\varepsilon_2$ at fixed $J$.
\end{enumerate}
\end{problem}

\begin{solution}
\textbf{(a)} By inspection, $\varepsilon_0 = \half(\varepsilon_1+\varepsilon_2)$,
$d_z = \varepsilon_1-\varepsilon_2$, $d_x = 2J$, $d_y = 0$.

\textbf{(b)} Identical chromophores means $d_z=0$, hence $\theta=\pi/2$ and the
eigenstates are $\ket{\pm} = \tfrac{1}{\sqrt2}\!\left(\ket{e_1g_2}\pm\ket{g_1e_2}\right)$
with $E_\pm = \varepsilon \pm J$. Note the sign: the symmetric combination is the
\emph{upper} state when $J>0$ and the \emph{lower} state when $J<0$.

\textbf{(c)} The oscillator strength goes as $|\bm{\mu}_1\pm\bm{\mu}_2|^2$. For two
parallel dipoles, the symmetric combination carries all of it and the antisymmetric
combination is dark. The point-dipole coupling is
$J = r^{-3}\!\left[\bm{\mu}_1\!\cdot\!\bm{\mu}_2 - 3(\bm{\mu}_1\!\cdot\!\hat{r})(\bm{\mu}_2\!\cdot\!\hat{r})\right]$:
side by side ($\bm{\mu}\perp\hat{r}$) gives $J=+\mu^2/r^3>0$, head to tail
($\bm{\mu}\parallel\hat{r}$) gives $J=-2\mu^2/r^3<0$.

\smallskip
\begin{center}\small
\begin{tabular}{@{}llll@{}}
\toprule
Geometry & Sign of $J$ & Bright (symmetric) state & Shift \\
\midrule
H-aggregate (side by side) & $J>0$ & upper, at $\varepsilon+J$ & \textbf{blue} \\
J-aggregate (head to tail) & $J<0$ & lower, at $\varepsilon+J$ & \textbf{red} \\
\bottomrule
\end{tabular}
\end{center}
\smallskip

\noindent
This is why J-aggregates show the narrow, red-shifted, intense band that makes them
useful as photographic sensitizers and light-harvesting mimics, while H-aggregates are
blue-shifted and comparatively dark.

\textbf{(d)} With $d_z\neq0$ the two chromophores are inequivalent, the states are no
longer exactly symmetric/antisymmetric, and the dark state steals intensity: the
oscillator strengths become $|\cos(\theta/2)\bm{\mu}_1 \pm \sin(\theta/2)\bm{\mu}_2|^2$
and neither state is completely dark. Energetic disorder therefore \emph{brightens}
H-aggregates --- a real and much-studied effect in molecular crystals.
\end{solution}

% ---------------------------------------------------------------------------
\begin{problem}{H\"uckel butadiene is two dimers}
\label{prob:butadiene}
The $\pi$ system of \emph{trans}-butadiene is a $4\times4$ H\"uckel problem
($H_{ii}=\alpha$, $H_{i,i+1}=\beta$), but it has a mirror symmetry exchanging
$1\leftrightarrow4$ and $2\leftrightarrow3$.
\begin{enumerate}[label=(\alph*),leftmargin=*,itemsep=1pt]
\item Form $s_{1,2}$ and $a_{1,2}$, the symmetric and antisymmetric combinations
$\tfrac{1}{\sqrt2}(p_1\pm p_4)$ and $\tfrac{1}{\sqrt2}(p_2\pm p_3)$, and show that
$\Hop$ block-diagonalizes into two $2\times2$ blocks.
\item Identify $d_z$ and $d_x$ for each block and use \cref{eq:eigenvalues} to obtain all
four $\pi$ energies. Check against the textbook result
$E = \alpha\pm1.618\beta,\ \alpha\pm0.618\beta$.
\item Verify by diagonalizing the full $4\times4$ matrix numerically.
\end{enumerate}
\end{problem}

\begin{solution}
\textbf{(a)} Sites 1 and 4 are not bonded, so $H_{14}=0$. Then
\[
\braket{s_1|\Hop|s_1} = \tfrac12(H_{11}+2H_{14}+H_{44}) = \alpha,
\qquad
\braket{s_2|\Hop|s_2} = \tfrac12(H_{22}+2H_{23}+H_{33}) = \alpha+\beta,
\]
\[
\braket{s_1|\Hop|s_2} = \tfrac12(H_{12}+H_{13}+H_{42}+H_{43}) = \beta,
\]
and every cross term $\braket{s_i|\Hop|a_j}$ vanishes because $\Hop$ commutes with the
mirror operation while $s$ and $a$ belong to different eigenvalues of it. Repeating for
the antisymmetric pair (the signs of $H_{23}$ and $H_{14}$ flip) gives
\[
\Hop_S = \begin{bmatrix}\alpha & \beta\\ \beta & \alpha+\beta\end{bmatrix},
\qquad
\Hop_A = \begin{bmatrix}\alpha & \beta\\ \beta & \alpha-\beta\end{bmatrix}.
\]

\textbf{(b)} Reading off the Pauli components:
\[
\begin{aligned}
\text{block } S:&\quad \varepsilon_0 = \alpha+\tfrac{\beta}{2}, &\quad d_z &= -\beta, &\quad d_x &= 2\beta;\\
\text{block } A:&\quad \varepsilon_0 = \alpha-\tfrac{\beta}{2}, &\quad d_z &= +\beta, &\quad d_x &= 2\beta.
\end{aligned}
\]
Both have $|\dvec| = \sqrt{\beta^2+4\beta^2} = \sqrt5\,|\beta|$, so
\[
E = \alpha \pm \frac{\beta}{2} \pm \frac{\sqrt5}{2}|\beta|
  \;=\; \alpha + 1.618\beta,\;\; \alpha+0.618\beta,\;\; \alpha-0.618\beta,\;\; \alpha-1.618\beta,
\]
since $\tfrac12(\sqrt5\pm1) = 1.618,\,0.618$ --- the golden ratio and its reciprocal,
which is where those famous H\"uckel numbers come from. With $\beta<0$ the two lowest
(occupied) levels are $\alpha+1.618\beta$ (from block $S$) and $\alpha+0.618\beta$
(from block $A$).

\textbf{(c)} Numerically, with $\alpha=0$ and $\beta=-1$: the full matrix gives
$\{-1.61803,\,-0.61803,\,+0.61803,\,+1.61803\}$, and the two blocks give
$\{-1.61803,\,+0.61803\}$ and $\{-0.61803,\,+1.61803\}$ respectively --- the same set.

\smallskip\noindent
\emph{Why this matters.} Reducing a large Hamiltonian to independent small blocks using
symmetry is one of the most useful manoeuvres in quantum chemistry, and here it turns a
$4\times4$ problem into two problems you can solve with \cref{eq:eigenvalues} in your
head. The same trick underlies symmetry-adapted linear combinations throughout group
theory.
\end{solution}

% ---------------------------------------------------------------------------
\begin{problem}{St\"uckelberg oscillations}
\label{prob:stuck}
Modify \texttt{landau\_zener\_numeric} so that instead of one linear sweep,
$d_z(t) = A\sin(\omega_{\rm sw}t)$ oscillates back and forth through the crossing.
\begin{enumerate}[label=(\alph*),leftmargin=*,itemsep=1pt]
\item Plot the population in $\ket{\alpha}$ against time over a few periods.
\item Explain why it does not simply relax toward $1/2$.
\item Find an amplitude and frequency for which the transfer after \emph{two} passages
is smaller than after one. What is interfering with what?
\end{enumerate}
\end{problem}

\begin{solution}
\textbf{(a),(b)} Each passage through the crossing acts as a beamsplitter, putting the
system into a superposition of the two adiabatic surfaces with amplitudes fixed by
\cref{eq:lz}. Between passages the two components accumulate \emph{different} dynamical
phases, $\Delta\phi = \hbar^{-1}\!\oint\!\left(E_+-E_-\right)\dd t$, because the surfaces
are separated. At the next passage the two amplitudes recombine and interfere. The
dynamics is therefore a repeated interferometer, not a repeated random draw --- so the
population does not diffuse to $1/2$, it oscillates. These are St\"uckelberg
oscillations, and the full sequence is a Landau--Zener--St\"uckelberg interferometer.

\textbf{(c)} Take $V_{DA}=0.35$, $A=6$ and scan $\omega_{\rm sw}$ with $\hbar=1$. The
ratio (transfer after two passages)/(transfer after one) is strongly non-monotonic ---
it runs from about $10$ down to about $1$ and back as $\omega_{\rm sw}$ varies over a
range of order unity. Two representative points:

\smallskip
\begin{center}\small
\begin{tabular}{@{}lccc@{}}
\toprule
$\omega_{\rm sw}$ & transfer after 1 passage & after 2 passages & ratio \\
\midrule
$0.450$ & $0.041$ & $0.430$ & $10.4$ \quad(constructive) \\
$0.700$ & $0.0244$ & $0.0241$ & $0.99$ \quad(destructive) \\
\bottomrule
\end{tabular}
\end{center}
\smallskip

\noindent
At $\omega_{\rm sw}\approx0.700$ the second passage almost exactly undoes the first: the
amplitude transferred on the way out and the amplitude transferred on the way back
arrive $\pi$ out of phase and cancel. Nothing has been ``untransferred'' in a classical
sense --- the two \emph{amplitudes} interfere destructively. Changing
$\omega_{\rm sw}$ by a few percent changes $\Delta\phi$ by order $\pi$ and flips the
result, which is why the ratio swings so violently.
\end{solution}

% ---------------------------------------------------------------------------
\begin{problem}{A genuine $\syo$: Peierls phases}
\label{prob:peierls}
A magnetic field attaches a phase to each hopping integral, $\beta\to\beta e^{\ii\phi_B}$.
\begin{enumerate}[label=(\alph*),leftmargin=*,itemsep=1pt]
\item Write $\Hop$ for a two-site dimer with this hopping and identify $d_x$ and $d_y$.
\item Show that the eigenvalues are independent of $\phi_B$, consistent with
\cref{sec:sigmay}.
\item Now take a triangle of three sites, each bond carrying phase $\phi_B$. Show that
the eigenvalues \emph{do} depend on the total flux $\Phi = 3\phi_B$ and that the
dependence cannot be rephased away. Diagonalize numerically and plot the three levels
against flux.
\end{enumerate}
\end{problem}

\begin{solution}
\textbf{(a)} With $H_{AB} = \beta e^{\ii\phi_B}$ and $\beta$ real,
$d_x = 2\beta\cos\phi_B$ and $d_y = -2\beta\sin\phi_B$, so
$\dperp = 2|\beta|$ regardless of $\phi_B$: the phase only rotates $\dvec$ within the
equatorial plane.

\textbf{(b)} By \cref{eq:eigenvalues} the eigenvalues depend on the transverse components
only through $\dperp$, so $E_\pm = \varepsilon_0 \pm |\beta|$ for every $\phi_B$.
Numerically, $\phi_B = 0,\,0.4,\,1.1,\,2.0$ all give $\{-|\beta|,+|\beta|\}$ exactly.
Equivalently: one rephasing, $\ket{B}\to e^{\ii\phi_B}\ket{B}$, removes the phase
completely.

\textbf{(c)} For the triangle with $t = \beta e^{\ii\phi_B}$ on each bond taken in a
consistent circulation,
\[
\Hop = \begin{bmatrix} 0 & t & t^{*}\\ t^{*} & 0 & t\\ t & t^{*} & 0\end{bmatrix},
\qquad
E_n = 2\beta\cos\!\left(\frac{2\pi n + \Phi}{3}\right),\quad n=0,1,2,
\qquad \Phi = 3\phi_B .
\]
This is verified numerically: at $\Phi=0$ the spectrum is
$\{-2|\beta|, +|\beta|, +|\beta|\}$ with the familiar doubly degenerate level of a
H\"uckel triangle; at $\Phi = \pi/2$ it is $\{-1.732,0,+1.732\}|\beta|$; at $\Phi = \pi$
it is $\{-|\beta|,-|\beta|,+2|\beta|\}$. The degeneracy is lifted and the whole spectrum
shifts with flux.

The counting argument is the point. Three sites admit three phase choices, but the
overall phase is unobservable, leaving \emph{two} useful rephasings against
\emph{three} bond phases. The combination that survives is the sum around the loop,
$\Phi = \phi_{12}+\phi_{23}+\phi_{31}$, which is gauge invariant and physically is the
magnetic flux threading the ring. A dimer has two sites and one bond --- one usable
rephasing, one phase --- so nothing survives. That is the whole difference between (b)
and (c), and it is the Aharonov--Bohm effect in its smallest possible setting.
\end{solution}

% ---------------------------------------------------------------------------
\begin{problem}{Ramsey interferometry}
\label{prob:ramsey}
Using the machinery of \cref{sec:driven}, simulate the following on a system with
detuning $\Delta$: a $\pi/2$ pulse with phase $\varphi=0$; free evolution with the drive
off (so $\dvec = \hbar\Delta\hat{z}$) for a time $\tau$; a second $\pi/2$ pulse, also
with $\varphi=0$.
\begin{enumerate}[label=(\alph*),leftmargin=*,itemsep=1pt]
\item Plot the final $P_\beta$ against $\tau$. What is the oscillation frequency?
\item Repeat with the second pulse phase-shifted to $\varphi=\pi$. What happens, and why
does this demonstrate that the \emph{relative} phase of two pulses is physical even
though the phase of one is not?
\item This is the Ramsey separated-oscillatory-fields method, and it is how atomic clocks
work. Why does a longer $\tau$ give a more precise frequency measurement?
\end{enumerate}
\end{problem}

\begin{solution}
Work in the hard-pulse limit $\omega_1\gg\Delta$, so the pulses are effectively
instantaneous rotations and the detuning acts only during $\tau$.

\textbf{(a)} The first $x$-pulse takes $\svec$ from the north pole to $-\hat{y}$
(\cref{sec:driven}). Free evolution rotates it in the equatorial plane at rate $\Delta$,
so $\svec(\tau) = (\sin\Delta\tau,\,-\cos\Delta\tau,\,0)$. The second $x$-pulse rotates
about $\hat{x}$ by $\ang{90}$ in the same sense ($\hat{y}\to\hat{z}$), leaving
$\svec = (\sin\Delta\tau,\,0,\,-\cos\Delta\tau)$. Hence
\[
P_\beta(\tau) = \frac{1-s_z}{2} = \frac{1+\cos\Delta\tau}{2},
\]
full-contrast fringes at the detuning frequency $\Delta$. (Numerically, with
$\omega_1 = 400\Delta$, the simulation matches this to better than $0.003$.)

\textbf{(b)} With $\varphi=\pi$ the second rotation is about $-\hat{x}$, so
$\hat{y}\to-\hat{z}$ and the final $s_z$ changes sign:
$P_\beta = \tfrac12(1-\cos\Delta\tau)$. The fringes are exactly inverted. Neither pulse
phase is meaningful alone --- calling the first one ``$x$'' was a convention --- but the
\emph{difference} of the two phases determines whether you measure a fringe maximum or a
minimum at $\tau=0$. That is a measurable, convention-independent statement, and it is
the concrete content of \cref{sec:sigmay}.

\textbf{(c)} The fringe pattern $\cos\Delta\tau$ has period $2\pi/\tau$ in $\Delta$, so a
frequency error $\delta\Delta$ produces a phase error $\delta\Delta\cdot\tau$. Doubling
$\tau$ halves the detuning uncertainty for the same phase resolution: the precision
improves as $1/\tau$. In practice $\tau$ is limited by decoherence, which is why
buying a better clock means buying a longer coherence time.
\end{solution}

% ---------------------------------------------------------------------------
\begin{problem}{The ammonia dipole paradox}
\label{prob:nh3}
At zero field the ground state of NH$_3$ is $\tfrac{1}{\sqrt2}(\ket{L}+\ket{R})$, which
has \emph{zero} expectation value of the dipole operator --- yet NH$_3$ is tabulated with
$\mu = \SI{1.47}{\debye}$.
\begin{enumerate}[label=(\alph*),leftmargin=*,itemsep=1pt]
\item Compute $\Braket{\hat\mu}$ in each of the two zero-field eigenstates, taking
$\hat\mu = \mu\,\szo$ in the $\{\ket{L},\ket{R}\}$ basis.
\item A molecule prepared in $\ket{L}$ is not stationary. Using \cref{eq:bloch-eq}, find
$\Braket{\hat\mu}(t)$ and its oscillation frequency.
\item Reconcile (a) and (b) with the tabulated value. On what timescale must a
measurement act to see $\SI{1.47}{\debye}$ rather than zero?
\item Deuterating to ND$_3$ drops the inversion splitting to about
\SI{0.053}{\wavenumber}. What happens to the crossover field
$\mathcal{E}_{\rm cross}=\Delta_t/\mu$, and why is the splitting so sensitive to
isotopic substitution?
\end{enumerate}
\end{problem}

\begin{solution}
\textbf{(a)} With $\hat\mu = \mu\szo$ and eigenstates
$\ket{\pm} = \tfrac{1}{\sqrt2}(\ket{L}\pm\ket{R})$,
\[
\Braket{\pm|\hat\mu|\pm} = \frac{\mu}{2}\big(\braket{L|\szo|L} + \braket{R|\szo|R}\big)
 = \frac{\mu}{2}(1-1) = 0 .
\]
Both stationary states are exactly nonpolar. This is not an accident: they are parity
eigenstates, and $\hat\mu$ is parity-odd, so its diagonal elements must vanish. The
dipole operator is purely \emph{off}-diagonal in the energy eigenbasis --- which is
precisely why the $\SI{23.786}{\giga\hertz}$ transition between them is dipole allowed
and drives the maser.

\textbf{(b)} $\ket{L} = \tfrac{1}{\sqrt2}(\ket{+}+\ket{-})$ has Bloch vector
$\svec(0) = +\hat{z}$, while $\dvec = -2\Delta_t\hat{x}$. By \cref{eq:bloch-eq} the
vector precesses in the $zy$ plane at $|\dvec|/\hbar = 2\Delta_t/\hbar$, so
\[
\Braket{\hat\mu}(t) = \mu\, s_z(t) = \mu\cos\!\left(\frac{2\Delta_t}{\hbar}t\right),
\]
oscillating at exactly the inversion frequency, \SI{23.786}{\giga\hertz}
(period \SI{42}{\pico\second}).

\textbf{(c)} Both are right; they answer different questions. A measurement slow
compared with \SI{42}{\pico\second} averages \cref{eq:bloch-eq} over many inversions and
returns zero. A measurement fast compared with it --- or any interaction that
\emph{localizes} the molecule before it can invert --- sees the instantaneous
\SI{1.47}{\debye}. In practice a bulk dielectric measurement reports
$\mu=\SI{1.47}{\debye}$ because collisions in the gas, and the local fields of
neighbouring molecules in a liquid, both act far faster than \SI{42}{\pico\second} and
keep the molecule pinned on one side. \Cref{fig:nh3} is the quantitative version of that
statement: a field of order $\mathcal{E}_{\rm cross}$ localizes the umbrella, and once
localized the molecule has its full classical dipole. The ``paradox'' is really the
statement that the tabulated dipole belongs to a \emph{localized} molecule, not to an
energy eigenstate of an isolated one.

\textbf{(d)} $\mathcal{E}_{\rm cross}$ is proportional to $\Delta_t$, so it falls by the
same factor of roughly $15$: from \SI{16.07}{\kilo\volt\per\centi\meter} to about
\SI{1.1}{\kilo\volt\per\centi\meter}. ND$_3$ is far easier to localize with a field.

The sensitivity comes from the exponential dependence of a tunneling amplitude on the
mass. In the WKB approximation
$\Delta_t \sim \exp\!\left(-\hbar^{-1}\!\int\!\sqrt{2m[V(x)-E]}\,\dd x\right)$, so
$\Delta_t$ falls roughly exponentially in $\sqrt{m}$. Deuterium is twice the mass of
hydrogen, and although the relevant inversion coordinate involves the nitrogen moving
against three hydrogens rather than a single H atom, the effective mass still increases
enough to suppress the amplitude by more than an order of magnitude. This exponential
mass sensitivity is the general signature of tunneling, and it is why kinetic isotope
effects far larger than the classical $\sqrt{2}$ are taken as evidence that a reaction
proceeds by tunneling.
\end{solution}


\appendix
% ============================================================================
\section{The Pauli algebra}
\label{app:pauli}
% ============================================================================

Everything in \cref{sec:formalism,sec:geometry} rests on one product rule. Multiplying
out the matrices directly,
%
\begin{equation}
\sxo\syo = \ii\szo,\qquad \syo\szo = \ii\sxo,\qquad \szo\sxo = \ii\syo,
\qquad
\sxo^2=\syo^2=\szo^2=\Iop,
\end{equation}
%
and each product in the first list reverses sign when the factors are exchanged. All of
this is summarized by

\begin{keyresult}[Pauli product rule]
\vspace{-0.35cm}
\begin{equation}
\hat\sigma_i\hat\sigma_j = \delta_{ij}\Iop + \ii\,\epsilon_{ijk}\hat\sigma_k ,
\label{eq:pauli-product}
\end{equation}
with $\epsilon_{ijk}$ the Levi-Civita symbol and a sum over $k$ implied.
\end{keyresult}

\noindent
Taking the symmetric and antisymmetric parts of \cref{eq:pauli-product} gives the
commutator and anticommutator separately,
%
\begin{equation}
\left[\hat\sigma_i,\hat\sigma_j\right] = 2\ii\,\epsilon_{ijk}\hat\sigma_k,
\qquad
\left\{\hat\sigma_i,\hat\sigma_j\right\} = 2\delta_{ij}\Iop .
\label{eq:pauli-comm}
\end{equation}
%
The first of these, written for $\Sxo=\tfrac{\hbar}{2}\sxo$ and friends, is the angular
momentum algebra $[\Sxo,\Syo]=\ii\hbar\Szo$ you verified numerically in Computational
Set 1. The trace identities \eqref{eq:traces} follow by taking the trace of
\cref{eq:pauli-product} and using $\Tr\hat\sigma_k=0$.

\subsection*{Proof of $(\dvec\cdot\pauliv)^2 = |\dvec|^2\Iop$}

Expand and apply \cref{eq:pauli-product}:
%
\begin{align}
\left(\dvec\cdot\pauliv\right)^2
 &= d_i d_j\,\hat\sigma_i\hat\sigma_j
  = d_i d_j\left(\delta_{ij}\Iop + \ii\,\epsilon_{ijk}\hat\sigma_k\right) \nonumber\\
 &= |\dvec|^2\,\Iop + \ii\,\epsilon_{ijk}\,d_i d_j\,\hat\sigma_k
  = |\dvec|^2\,\Iop .
\label{eq:dsq-proof}
\end{align}
%
The second term vanishes because $\epsilon_{ijk}$ is antisymmetric in $i\leftrightarrow j$
while $d_i d_j$ is symmetric, and the contraction of a symmetric with an antisymmetric
tensor is zero. This is \cref{eq:dsquared}, and with it the eigenvalues
\eqref{eq:eigenvalues} follow immediately: if $\hat{M}^2 = |\dvec|^2\Iop$ then every
eigenvalue $m$ of $\hat{M}$ satisfies $m^2=|\dvec|^2$.

% ============================================================================
\section{The eigenvectors}
\label{app:eigvec}
% ============================================================================

Take $\varepsilon_0=0$ without loss of generality and write $\Hop$ in terms of the polar
angles of \cref{eq:mixing-angle}:
%
\begin{equation}
\Hop = \frac{|\dvec|}{2}
\begin{bmatrix}
\cos\theta & \sin\theta\,e^{-\ii\varphi}\\
\sin\theta\,e^{+\ii\varphi} & -\cos\theta
\end{bmatrix}.
\end{equation}
%
Try $\ket{+} = \big(\cos\tfrac{\theta}{2},\; e^{\ii\varphi}\sin\tfrac{\theta}{2}\big)^{T}$.
The first component of $\Hop\ket{+}$ is
%
\begin{align}
\frac{|\dvec|}{2}\left[\cos\theta\cos\tfrac{\theta}{2}
 + \sin\theta\,e^{-\ii\varphi}e^{\ii\varphi}\sin\tfrac{\theta}{2}\right]
&= \frac{|\dvec|}{2}\left[\cos\theta\cos\tfrac{\theta}{2} + \sin\theta\sin\tfrac{\theta}{2}\right] \nonumber\\
&= \frac{|\dvec|}{2}\cos\!\left(\theta-\tfrac{\theta}{2}\right)
 = \frac{|\dvec|}{2}\cos\tfrac{\theta}{2},
\end{align}
%
using the cosine difference formula. The second component is, similarly,
%
\begin{equation}
\frac{|\dvec|}{2}\,e^{\ii\varphi}\left[\sin\theta\cos\tfrac{\theta}{2} - \cos\theta\sin\tfrac{\theta}{2}\right]
= \frac{|\dvec|}{2}\,e^{\ii\varphi}\sin\tfrac{\theta}{2}.
\end{equation}
%
So $\Hop\ket{+} = \tfrac{|\dvec|}{2}\ket{+}$, confirming \cref{eq:eigenvectors} and
$E_+ = \varepsilon_0 + |\dvec|/2$. The state $\ket{-}$ is obtained by
$\theta\to\theta+\pi$, $\varphi\to\varphi+\pi$ (that is, $\dvec\to-\dvec$), which is why
it appears with the opposite sign pattern.

Two checks worth doing yourself: $\braket{+|-} = 0$, and
$\braket{+|\szo|+} = \cos^2\tfrac{\theta}{2}-\sin^2\tfrac{\theta}{2} = \cos\theta$, which
together with the analogous $x$ and $y$ expectation values confirms
$\svec_+ = \dhat$ as asserted in \cref{sec:bloch}.

% ============================================================================
\section{The Bloch equation}
\label{app:blocheq}
% ============================================================================

Start from the Heisenberg equation for the expectation value of $\hat\sigma_a$,
%
\begin{equation}
\frac{\dd\Braket{\hat\sigma_a}}{\dd t}
  = \frac{\ii}{\hbar}\Braket{\left[\Hop,\hat\sigma_a\right]}.
\end{equation}
%
The identity part of $\Hop$ commutes with everything, so only
$\half d_b\hat\sigma_b$ contributes:
%
\begin{equation}
\left[\Hop,\hat\sigma_a\right]
 = \half d_b\left[\hat\sigma_b,\hat\sigma_a\right]
 = \half d_b\left(2\ii\,\epsilon_{bac}\hat\sigma_c\right)
 = \ii\, d_b\,\epsilon_{bac}\,\hat\sigma_c ,
\end{equation}
%
using \cref{eq:pauli-comm}. Therefore
%
\begin{equation}
\frac{\dd\Braket{\hat\sigma_a}}{\dd t}
 = \frac{\ii}{\hbar}\cdot \ii\,d_b\,\epsilon_{bac}\Braket{\hat\sigma_c}
 = -\frac{1}{\hbar}\,\epsilon_{bac}\,d_b\Braket{\hat\sigma_c}
 = +\frac{1}{\hbar}\,\epsilon_{abc}\,d_b\Braket{\hat\sigma_c},
\end{equation}
%
where the last step used $\epsilon_{bac} = -\epsilon_{abc}$. The right-hand side is
precisely the $a$-component of a cross product,
$\left(\dvec\times\svec\right)_a = \epsilon_{abc}d_b s_c$, giving \cref{eq:bloch-eq}:
%
\begin{equation}
\frac{\dd\svec}{\dd t} = \frac{1}{\hbar}\,\dvec\times\svec .
\end{equation}

\noindent
\textbf{Remark.} For a real spin in a real magnetic field, $\dvec = -\gamma\hbar\mathbf{B}$
and this becomes $\dot{\svec} = -\gamma\,\mathbf{B}\times\svec$ --- the classical
equation for a magnetic moment precessing in a field, and the reason the same picture
works for a classical gyroscope. That a fully quantum two-level system obeys a
\emph{closed, classical-looking} equation of motion is special to two dimensions: the
Bloch vector contains all the information in the state, so no higher moments are needed
to close the hierarchy.

% ============================================================================
\section{The rotating-frame transformation}
\label{app:rotframe}
% ============================================================================

Let $\Ry(t) = e^{+\ii\omega t\szo/2}$ and $\ket{\tilde\psi} = \Ry\ket{\psi}$. Then
%
\begin{align}
\ii\hbar\frac{\dd\ket{\tilde\psi}}{\dd t}
 &= \ii\hbar\dot{\Ry}\ket{\psi} + \Ry\left(\ii\hbar\frac{\dd\ket{\psi}}{\dd t}\right)
  = \ii\hbar\left(\frac{\ii\omega}{2}\szo\right)\Ry\ket{\psi} + \Ry\Hop\ket{\psi} \nonumber\\
 &= \left[\Ry\Hop\Ry^{\dagger} - \frac{\hbar\omega}{2}\szo\right]\ket{\tilde\psi}
  \;\equiv\; \Hop_{\rm rot}\ket{\tilde\psi}.
\label{eq:rotframe-def}
\end{align}
%
So we need $\Ry\Hop\Ry^\dagger$. Since $\Ry$ is built from $\szo$, the $\szo$ term passes
through untouched. For the transverse operators, define
$\hat{A}(\vartheta) = e^{\ii\vartheta\szo/2}\,\sxo\,e^{-\ii\vartheta\szo/2}$ and
differentiate:
%
\begin{equation}
\frac{\dd\hat{A}}{\dd\vartheta} = \frac{\ii}{2}\left[\szo,\hat{A}\right],
\qquad
\left.\frac{\dd\hat{A}}{\dd\vartheta}\right|_{0} = \frac{\ii}{2}\left(2\ii\syo\right) = -\syo .
\end{equation}
%
The analogous calculation for $\hat{B}(\vartheta)$ built from $\syo$ gives
$\hat{B}'(0) = +\sxo$, so the pair satisfies $\hat{A}'=-\hat{B}$, $\hat{B}'=+\hat{A}$ and
%
\begin{equation}
\Ry\,\sxo\,\Ry^{\dagger} = \sxo\cos\omega t - \syo\sin\omega t,
\qquad
\Ry\,\syo\,\Ry^{\dagger} = \syo\cos\omega t + \sxo\sin\omega t.
\label{eq:rot-transverse}
\end{equation}
%
Now substitute \cref{eq:rot-transverse} into the drive term of \cref{eq:lab-frame}. Write
$c=\cos(\omega t+\varphi)$, $s=\sin(\omega t+\varphi)$, $C=\cos\omega t$,
$S=\sin\omega t$:
%
\begin{align}
\frac{\hbar\omega_1}{2}\Big[c\left(\sxo C - \syo S\right) + s\left(\syo C + \sxo S\right)\Big]
&= \frac{\hbar\omega_1}{2}\Big[\sxo\left(cC+sS\right) + \syo\left(sC-cS\right)\Big] \nonumber\\
&= \frac{\hbar\omega_1}{2}\Big[\sxo\cos\varphi + \syo\sin\varphi\Big],
\end{align}
%
where the two brackets collapsed by the cosine and sine difference formulas evaluated at
$(\omega t+\varphi)-\omega t = \varphi$. \textbf{All the time dependence has cancelled.}
Combining with the $\szo$ terms from \cref{eq:rotframe-def},
%
\begin{equation}
\Hop_{\rm rot} = \frac{\hbar\omega_0}{2}\szo - \frac{\hbar\omega}{2}\szo
 + \frac{\hbar\omega_1}{2}\left(\cos\varphi\,\sxo+\sin\varphi\,\syo\right),
\end{equation}
%
which is \cref{eq:rotframe}.

\noindent
\textbf{What was given up.} Nothing was approximated --- for a genuinely
\emph{circularly} rotating drive this transformation is exact. (For a linearly polarized
drive, $\propto\cos\omega t\,\sxo$, one first decomposes into two counter-rotating
components and discards the far-off-resonant one; that discard is the rotating-wave
approximation, and it is a genuine approximation.) What is lost is directness of
interpretation: $\ket{\tilde\psi}$ is not the physical state, and a stationary Bloch
vector in the rotating frame corresponds to one spinning at $\omega$ in the laboratory.

% ============================================================================
\section{The Marcus barrier}
\label{app:marcus}
% ============================================================================

From \cref{eq:diabats} with $\lambda = 2kq_0^2$, the diabats cross where
$V_D(q^\ddagger)=V_A(q^\ddagger)$:
%
\begin{equation}
\half k(q+q_0)^2 = \half k(q-q_0)^2 + \Delta G^\circ
\;\Longrightarrow\;
2kq_0 q = \Delta G^\circ
\;\Longrightarrow\;
q^{\ddagger} = \frac{\Delta G^\circ}{2kq_0} = \frac{\Delta G^\circ q_0}{\lambda}.
\end{equation}
%
The activation energy is the diabatic energy there, measured from the donor minimum at
$q=-q_0$ where $V_D=0$:
%
\begin{align}
\Delta G^{\ddagger} = V_D(q^{\ddagger})
 &= \half k\left(\frac{\Delta G^\circ q_0}{\lambda} + q_0\right)^{\!2}
  = \half k q_0^2\left(\frac{\Delta G^\circ+\lambda}{\lambda}\right)^{\!2} \nonumber\\
 &= \frac{\lambda}{4}\cdot\frac{(\lambda+\Delta G^\circ)^2}{\lambda^2}
  = \frac{(\lambda+\Delta G^{\circ})^2}{4\lambda},
\end{align}
%
using $\half k q_0^2 = \lambda/4$. This is \cref{eq:marcus-barrier}. Three limits are
worth memorizing:
%
\begin{equation}
\Delta G^\circ = 0 \;\Rightarrow\; \Delta G^\ddagger = \frac{\lambda}{4},
\qquad
\Delta G^\circ = -\lambda \;\Rightarrow\; \Delta G^\ddagger = 0,
\qquad
\Delta G^\circ = -2\lambda \;\Rightarrow\; \Delta G^\ddagger = \frac{\lambda}{4}.
\end{equation}
%
The barrier is symmetric about $\Delta G^\circ=-\lambda$: a reaction that is
\emph{too} exergonic by exactly $\lambda$ is as slow as one with no driving force at all.

Note finally that this is the barrier on the \emph{diabatic} surface. The true barrier on
the lower adiabatic surface is lower by roughly $V_{DA}$ near the crossing. When
$V_{DA}\ll k_BT$ the nonadiabatic (golden-rule) rate applies and the diabatic barrier is
the right one to use; when $V_{DA}\gg k_BT$ the reaction is adiabatic and the transition
state is the top of $E_-(q)$. The dimensionless ratio that decides which regime you are
in is essentially the Landau--Zener exponent of \cref{eq:lz}.


% ============================================================================
%  References. Kept as a thebibliography environment so the project compiles in
%  one pdfLaTeX pass with no BibTeX/Biber run. To switch to a .bib workflow,
%  delete this file, add \usepackage[style=chem-acs]{biblatex} +
%  \addbibresource{refs.bib} to preamble.sty, and \printbibliography here.
% ============================================================================
\begin{thebibliography}{99}
\small

\bibitem{Huckel1931}
E.~H\"uckel, ``Quantentheoretische Beitr\"age zum Benzolproblem,''
\emph{Z. Phys.} \textbf{70}, 204--286 (1931).

\bibitem{Lowdin1950}
P.-O.~L\"owdin, ``On the Non-Orthogonality Problem Connected with the Use of Atomic
Wave Functions in the Theory of Molecules and Crystals,''
\emph{J. Chem. Phys.} \textbf{18}, 365--375 (1950).

\bibitem{RobinDay}
M.~B.~Robin and P.~Day, ``Mixed Valence Chemistry --- A Survey and Classification,''
\emph{Adv. Inorg. Chem. Radiochem.} \textbf{10}, 247--422 (1968).

\bibitem{CreutzTaube}
C.~Creutz and H.~Taube, ``A Direct Approach to Measuring the Franck--Condon Barrier to
Electron Transfer between Metal Ions,''
\emph{J. Am. Chem. Soc.} \textbf{91}, 3988--3989 (1969).

\bibitem{Rabi1937}
I.~I.~Rabi, ``Space Quantization in a Gyrating Magnetic Field,''
\emph{Phys. Rev.} \textbf{51}, 652--654 (1937).

\bibitem{Feynman1957}
R.~P.~Feynman, F.~L.~Vernon and R.~W.~Hellwarth, ``Geometrical Representation of the
Schr\"odinger Equation for Solving Maser Problems,''
\emph{J. Appl. Phys.} \textbf{28}, 49--52 (1957).
\emph{(Introduced the Bloch-vector picture for a general two-level system, not just
spins --- the paper that justifies this entire set of notes.)}

\bibitem{Ramsey1950}
N.~F.~Ramsey, ``A Molecular Beam Resonance Method with Separated Oscillating Fields,''
\emph{Phys. Rev.} \textbf{78}, 695--699 (1950).

\bibitem{Marcus1956}
R.~A.~Marcus, ``On the Theory of Oxidation--Reduction Reactions Involving Electron
Transfer. I,'' \emph{J. Chem. Phys.} \textbf{24}, 966--978 (1956).

\bibitem{MarcusNobel}
R.~A.~Marcus, ``Electron Transfer Reactions in Chemistry: Theory and Experiment
(Nobel Lecture),'' \emph{Rev. Mod. Phys.} \textbf{65}, 599--610 (1993).

\bibitem{Landau1932}
L.~D.~Landau, ``Zur Theorie der Energie\"ubertragung II,''
\emph{Phys. Z. Sowjetunion} \textbf{2}, 46--51 (1932).

\bibitem{Zener1932}
C.~Zener, ``Non-Adiabatic Crossing of Energy Levels,''
\emph{Proc. R. Soc. London A} \textbf{137}, 696--702 (1932).

\bibitem{Stuckelberg1932}
E.~C.~G.~St\"uckelberg, ``Theorie der unelastischen St\"osse zwischen Atomen,''
\emph{Helv. Phys. Acta} \textbf{5}, 369--422 (1932).

\bibitem{Townes1955}
J.~P.~Gordon, H.~J.~Zeiger and C.~H.~Townes, ``The Maser --- New Type of Microwave
Amplifier, Frequency Standard, and Spectrometer,''
\emph{Phys. Rev.} \textbf{99}, 1264--1274 (1955).

\bibitem{JaynesCummings}
E.~T.~Jaynes and F.~W.~Cummings, ``Comparison of Quantum and Semiclassical Radiation
Theories with Application to the Beam Maser,''
\emph{Proc. IEEE} \textbf{51}, 89--109 (1963).

\bibitem{Holstein1959}
T.~Holstein, ``Studies of Polaron Motion: Part I. The Molecular-Crystal Model,''
\emph{Ann. Phys.} \textbf{8}, 325--342 (1959).

\bibitem{Leggett1987}
A.~J.~Leggett, S.~Chakravarty, A.~T.~Dorsey, M.~P.~A.~Fisher, A.~Garg and W.~Zwerger,
``Dynamics of the Dissipative Two-State System,''
\emph{Rev. Mod. Phys.} \textbf{59}, 1--85 (1987).

\bibitem{Kasha1965}
M.~Kasha, H.~R.~Rawls and M.~A.~El-Bayoumi, ``The Exciton Model in Molecular
Spectroscopy,'' \emph{Pure Appl. Chem.} \textbf{11}, 371--392 (1965).
\emph{(H- and J-aggregates; Problem~\ref{prob:exciton}.)}

\bibitem{ShevchenkoLZS}
S.~N.~Shevchenko, S.~Ashhab and F.~Nori, ``Landau--Zener--St\"uckelberg
Interferometry,'' \emph{Phys. Rep.} \textbf{492}, 1--30 (2010).
\emph{(Problem~\ref{prob:stuck}.)}

\end{thebibliography}


\end{document}
